% Requires in the preamble:
% \usepackage{amsmath,amssymb}
% \usepackage{tikz,pgfplots}
% \pgfplotsset{compat=1.18}

\chapter{The Integral as Accumulation}
\label{chap:integral-as-accumulation}

The last chapter showed how much a derivative can do.  A derivative can relate
changing quantities, find a best choice, estimate error, and express laws such as

\[
T'=-k(T-T_s)
\]

or

\[
P'=rP\left(1-\frac{P}{K}\right).
\]

But a derivative is local.  It tells what is happening at an instant.

Many questions ask for the accumulated result of those instant-by-instant changes.

\[
\text{If velocity is known, how far did the object travel?}
\]

\[
\text{If power is used over time, how much energy was used?}
\]

\[
\text{If density varies along a rod, what is the rod's total mass?}
\]

\[
\text{If a rate of change is known, what is the total change?}
\]

These questions return us to the dashboard from the beginning of the book.  The
speedometer gives a rate.  The odometer gives an accumulated amount.

Differentiation went from accumulation to rate:

\[
\text{position} \longrightarrow \text{velocity}.
\]

Integration goes in the opposite direction:

\[
\text{velocity} \longrightarrow \text{position change}.
\]

This chapter builds the integral from the ground up.  We will not begin with
antiderivatives.  We will begin with the simpler and older idea: break a quantity into
small pieces, estimate each piece, and add.

\section{Distance from velocity}
\label{sec:distance-from-velocity}

A derivative tells how fast a quantity is changing.  To recover the quantity itself, we
must add up many small changes.

If \(s(t)\) is position and \(v(t)=s'(t)\) is velocity, then for a short time interval of
length \(\Delta t\),

\[
\Delta s \approx v(t)\Delta t.
\]

This is the local linear idea in reverse.  The derivative says

\[
s(t+\Delta t)\approx s(t)+v(t)\Delta t.
\]

So the small position change is approximately

\[
v(t)\Delta t.
\]

The integral begins by adding many expressions of this kind.

\subsection{Constant velocity revisited}
\label{subsec:constant-velocity-revisited}

Suppose a car moves at a constant velocity

\[
v(t)=60
\]

miles per hour from \(t=0\) to \(t=3\).  Then its displacement is

\[
60\cdot 3=180
\]

miles.

There is no calculus hidden here.  It is only

\[
\text{distance}=\text{velocity}\cdot\text{time}.
\]

But the graph already contains the key idea.  The velocity graph is horizontal.  The
region under it from \(0\) to \(3\) is a rectangle with height \(60\) and width \(3\).

\[
\text{rectangle area}=60\cdot 3=180.
\]

The units also agree:

\[
\frac{\text{miles}}{\text{hour}}\cdot \text{hour}=\text{miles}.
\]

\begin{figure}[h]
\centering
\begin{tikzpicture}
\begin{axis}[
    width=0.78\textwidth,
    height=0.45\textwidth,
    xmin=0, xmax=3.4,
    ymin=0, ymax=75,
    axis lines=left,
    xlabel={\(t\) (hours)},
    ylabel={\(v(t)\) (miles/hour)},
    xtick={0,1,2,3},
    ytick={0,60},
    grid=both,
    grid style={line width=.1pt, draw=gray!25},
]
\draw[fill=gray!15, draw=gray!60] (axis cs:0,0) rectangle (axis cs:3,60);
\addplot[very thick, domain=0:3, samples=2] {60};
\node at (axis cs:1.5,30) {\(\text{area}=60\cdot3\)};
\node[anchor=west] at (axis cs:3.05,60) {\(v=60\)};
\end{axis}
\end{tikzpicture}
\caption{For constant velocity, displacement is the area of a rectangle under the velocity graph.}
\label{fig:constant-velocity-area}
\end{figure}

If the same constant velocity lasts from \(t=a\) to \(t=b\), then the displacement is

\[
60(b-a).
\]

More generally, if

\[
v(t)=v_0
\]

on the interval \([a,b]\), then

\[
\boxed{
\text{displacement}=v_0(b-a).
}
\]

If the initial position is \(s(a)\), then

\[
\boxed{
s(b)=s(a)+v_0(b-a).
}
\]

This is the simplest accumulation formula.

\paragraph{Example.}
A train is already \(12\) miles from a station at noon.  It moves away from the station
at \(50\) miles per hour until \(2{:}30\) p.m.  What is its position then?

The time elapsed is

\[
2.5\text{ hours}.
\]

The displacement is

\[
50(2.5)=125
\]

miles.  Since the train started \(12\) miles from the station,

\[
s(2.5)=12+125=137.
\]

So the train is \(137\) miles from the station at \(2{:}30\) p.m.

The starting value \(12\) is not produced by the velocity.  The velocity produces the
change.  Accumulation adds that change to the starting value.

\subsection{Piecewise constant velocity}
\label{subsec:piecewise-constant-velocity}

Real velocity is rarely constant for a whole trip.  The next simplest case is
\emph{piecewise constant} velocity: constant for a while, then another constant, then
another.

Suppose a cart moves with the following velocity.

\[
v(t)=
\begin{cases}
30, & 0\le t<1,\\
50, & 1\le t<3,\\
20, & 3\le t\le 4,
\end{cases}
\]

where time is measured in hours and velocity in miles per hour.

During the first hour, the cart travels

\[
30\cdot1=30
\]

miles.

During the next two hours, it travels

\[
50\cdot2=100
\]

miles.

During the last hour, it travels

\[
20\cdot1=20
\]

miles.

So the total displacement is

\[
30+100+20=150
\]

miles.

The graph tells the same story.  The area under the velocity graph is the sum of three
rectangle areas.

\begin{figure}[h]
\centering
\begin{tikzpicture}
\begin{axis}[
    width=0.82\textwidth,
    height=0.48\textwidth,
    xmin=0, xmax=4.3,
    ymin=0, ymax=60,
    axis lines=left,
    xlabel={\(t\) (hours)},
    ylabel={\(v(t)\) (miles/hour)},
    xtick={0,1,2,3,4},
    ytick={0,20,30,50},
    grid=both,
    grid style={line width=.1pt, draw=gray!25},
]
\draw[fill=gray!12, draw=gray!60] (axis cs:0,0) rectangle (axis cs:1,30);
\draw[fill=gray!22, draw=gray!60] (axis cs:1,0) rectangle (axis cs:3,50);
\draw[fill=gray!12, draw=gray!60] (axis cs:3,0) rectangle (axis cs:4,20);

\addplot[very thick] coordinates {(0,30) (1,30)};
\addplot[very thick] coordinates {(1,50) (3,50)};
\addplot[very thick] coordinates {(3,20) (4,20)};

\node at (axis cs:0.5,15) {\(30\cdot1\)};
\node at (axis cs:2,25) {\(50\cdot2\)};
\node at (axis cs:3.5,10) {\(20\cdot1\)};
\end{axis}
\end{tikzpicture}
\caption{For piecewise constant velocity, displacement is a sum of rectangle areas.}
\label{fig:piecewise-constant-velocity}
\end{figure}

This is already the basic structure of integration:

\[
\text{total change}
=
\text{first small change}
+
\text{second small change}
+
\cdots.
\]

For piecewise constant velocity, the small changes are exact.  On each interval,

\[
\text{change in position}
=
\text{velocity on that interval}\cdot\text{length of that interval}.
\]

\paragraph{Example.}
A cyclist's velocity is

\[
v(t)=
\begin{cases}
12, & 0\le t<0.5,\\
18, & 0.5\le t<1.5,\\
10, & 1.5\le t\le 2,
\end{cases}
\]

where \(t\) is measured in hours and \(v\) in miles per hour.  Find the displacement
from \(t=0\) to \(t=2\).

The three time intervals have lengths

\[
0.5,\qquad 1,\qquad 0.5.
\]

So the displacement is

\[
12(0.5)+18(1)+10(0.5)=6+18+5=29.
\]

The cyclist travels \(29\) miles.

\subsection{Estimating distance from sampled speeds}
\label{subsec:estimating-distance-from-sampled-speeds}

In practice, we often do not know a velocity formula.  We may only have speedometer
readings at certain times.

Suppose a car's velocity is sampled every half hour.

\[
\begin{array}{c|ccccc}
t\text{ (hours)} & 0 & 0.5 & 1.0 & 1.5 & 2.0 \\ \hline
v(t)\text{ (miles/hour)} & 20 & 28 & 34 & 38 & 40
\end{array}
\]

How far did the car travel in the two hours?

The table does not tell us exactly what happened between readings.  But we can
estimate.

If we use the velocity at the beginning of each half-hour interval, we get

\[
0.5(20)+0.5(28)+0.5(34)+0.5(38).
\]

Thus

\[
\text{left estimate}
=
0.5(20+28+34+38)
=
60
\]

miles.

If we use the velocity at the end of each half-hour interval, we get

\[
0.5(28)+0.5(34)+0.5(38)+0.5(40).
\]

Thus

\[
\text{right estimate}
=
0.5(28+34+38+40)
=
70
\]

miles.

The two estimates are different because the velocity is changing.  If the velocity was
increasing throughout the trip, the left estimate is too small and the right estimate is
too large.

\begin{quote}
\textbf{The principle of rectangular estimation.}
On a short interval, pretend the velocity is constant.  Then

\[
\text{small displacement}\approx \text{sampled velocity}\cdot \text{time width}.
\]

Add the small displacements.
\end{quote}

If we also know midpoint speeds, we can often do better.  Suppose the readings at the
midpoints of the four half-hour intervals are

\[
\begin{array}{c|cccc}
t\text{ (hours)} & 0.25 & 0.75 & 1.25 & 1.75 \\ \hline
v(t)\text{ (miles/hour)} & 24 & 31 & 36 & 39
\end{array}
\]

Then the midpoint estimate is

\[
0.5(24+31+36+39)=65
\]

miles.

The estimates

\[
60,\qquad 65,\qquad 70
\]

are not three different truths.  They are three attempts to recover one accumulated
distance from incomplete velocity information.

Smaller time intervals usually give better estimates.  If we sample every minute
instead of every half hour, each rectangle covers less time, so the assumption
``velocity is nearly constant here'' becomes more trustworthy.

\begin{figure}[h]
\centering
\begin{tikzpicture}
\begin{axis}[
    width=0.82\textwidth,
    height=0.6\textwidth,
    xmin=0, xmax=2.05,
    ymin=0, ymax=45,
    axis lines=left,
    xlabel={\(t\) (hours)},
    ylabel={\(v(t)\) (miles/hour)},
    xtick={0,0.5,1,1.5,2},
    ytick={0,20,28,34,38,40},
    grid=both,
    grid style={line width=.1pt, draw=gray!25},
]
% left rectangles
\draw[fill=gray!12, draw=gray!55] (axis cs:0,0) rectangle (axis cs:0.5,20);
\draw[fill=gray!12, draw=gray!55] (axis cs:0.5,0) rectangle (axis cs:1,28);
\draw[fill=gray!12, draw=gray!55] (axis cs:1,0) rectangle (axis cs:1.5,34);
\draw[fill=gray!12, draw=gray!55] (axis cs:1.5,0) rectangle (axis cs:2,38);

% a smooth suggested velocity curve
\addplot[very thick, domain=0:2, samples=100] {-4*x^2 + 18*x + 20};

% measured points
\addplot[only marks, mark=*] coordinates {(0,20) (0.5,28) (1,34) (1.5,38) (2,40)};

\node[anchor=west] at (axis cs:0.1,42) {left rectangles from sampled velocities};
\end{axis}
\end{tikzpicture}
\caption{A sampled velocity table leads naturally to rectangular distance estimates.}
\label{fig:sampled-velocity-rectangles}
\end{figure}

This is the practical heart of the integral.  We do not need a perfect formula at the
start.  We need a way to add many small rate-times-width contributions.

\subsection{Signed displacement versus total distance}
\label{subsec:signed-displacement-total-distance}

Velocity has a sign.  Positive velocity means motion in the positive direction.
Negative velocity means motion in the negative direction.

That creates two different questions.

\[
\text{What is the net change in position?}
\]

\[
\text{How much distance was actually traveled?}
\]

The first question asks for \emph{signed displacement}.  The second asks for
\emph{total distance traveled}.

\paragraph{Example.}
A particle moves along a line with velocity

\[
v(t)=
\begin{cases}
4, & 0\le t<3,\\
-2, & 3\le t\le 8,
\end{cases}
\]

where \(t\) is measured in seconds and \(v(t)\) in feet per second.

From \(0\) to \(3\), the particle moves

\[
4\cdot3=12
\]

feet in the positive direction.

From \(3\) to \(8\), the particle moves

\[
(-2)\cdot5=-10
\]

feet in the negative direction.

So the signed displacement is

\[
12+(-10)=2
\]

feet.

The particle ends \(2\) feet to the right of where it started.

But the total distance traveled is

\[
12+10=22
\]

feet.

The particle traveled \(12\) feet one way and \(10\) feet back.

\begin{figure}[h]
\centering
\begin{tikzpicture}
\begin{axis}[
    width=0.82\textwidth,
    height=0.6\textwidth,
    xmin=0, xmax=9.0,
    ymin=-3, ymax=5,
    axis lines=middle,
    xlabel={\(t\) (seconds)},
    ylabel={\(v(t)\) (feet/second)},
    xtick={0,3,8},
    ytick={-2,0,4},
    grid=both,
    grid style={line width=.1pt, draw=gray!25},
]
\draw[fill=gray!18, draw=gray!60, opacity=0.5] (axis cs:0,0) rectangle (axis cs:3,4);
\draw[fill=gray!35, draw=gray!60, opacity=0.5] (axis cs:3,0) rectangle (axis cs:8,-2);

\addplot[very thick] coordinates {(0,4) (3,4)};
\addplot[very thick] coordinates {(3,-2) (8,-2)};

\node at (axis cs:1.5,2) {\(+12\)};
\node at (axis cs:5.5,-1) {\(-10\)};
\node[anchor=west] at (axis cs:0.2,4.25) {positive displacement};
\node[anchor=west] at (axis cs:3.2,-2.5) {negative displacement};
\end{axis}
\end{tikzpicture}
\caption{Area above the time axis counts positive.  Area below the time axis counts negative for displacement.}
\label{fig:signed-velocity-area}
\end{figure}

For piecewise constant velocity,

\[
\text{signed displacement}
=
v_1\Delta t_1+v_2\Delta t_2+\cdots+v_n\Delta t_n.
\]

The total distance traveled is

\[
\text{total distance}
=
|v_1|\Delta t_1+|v_2|\Delta t_2+\cdots+|v_n|\Delta t_n.
\]

The absolute values remove the signs because distance traveled does not cancel.

\begin{quote}
\textbf{Two accumulations from velocity.}

Signed displacement accumulates velocity:

\[
\text{displacement}=\text{positive contributions}+\text{negative contributions}.
\]

Total distance accumulates speed:

\[
\text{distance traveled}=\text{accumulation of } |v(t)|.
\]
\end{quote}

The picture suggests a new geometric language.  Positive velocity gives area above
the axis.  Negative velocity gives area below the axis.  For displacement, area below
the axis counts negatively.

So the next problem is geometric:

\[
\text{How do we measure area under a curved graph?}
\]

\section{Area under a graph}
\label{sec:area-under-a-graph}

When velocity is nonnegative, distance traveled is the area under the velocity graph.
That observation is too useful to keep only for motion.

Given a nonnegative function \(f\) on an interval \([a,b]\), we want to measure the area
between the graph and the \(x\)-axis.

For a rectangle, area is easy:

\[
\text{area}=\text{height}\cdot\text{width}.
\]

For a curved graph, the height changes.  The central idea is to replace the curved
region by many thin rectangles.

\[
\text{curved area}
\approx
\text{sum of many rectangular areas}.
\]

As the rectangles become thinner, the approximation should improve.

\subsection{Rectangular estimates}
\label{subsec:rectangular-estimates}

Consider the area under

\[
f(x)=x^2
\]

from \(x=0\) to \(x=2\).

The region is not a triangle or a rectangle.  We do not yet have an exact formula for
its area.  But we can estimate it.

Divide \([0,2]\) into four equal intervals:

\[
[0,0.5],\qquad [0.5,1],\qquad [1,1.5],\qquad [1.5,2].
\]

Each interval has width

\[
\Delta x=0.5.
\]

On each interval, choose one height from the function and build a rectangle.

If we use the left endpoint of each interval, the heights are

\[
f(0),\qquad f(0.5),\qquad f(1),\qquad f(1.5).
\]

Thus the rectangular estimate is

\[
0.5f(0)+0.5f(0.5)+0.5f(1)+0.5f(1.5).
\]

Since \(f(x)=x^2\),

\[
0.5\left(0^2+0.5^2+1^2+1.5^2\right)
=
0.5(0+0.25+1+2.25).
\]

Therefore

\[
\text{left rectangle estimate}=1.75.
\]

\begin{figure}[h]
\centering
\begin{tikzpicture}
\begin{axis}[
    width=0.78\textwidth,
    height=0.48\textwidth,
    xmin=0, xmax=2.2,
    ymin=0, ymax=4.4,
    axis lines=left,
    xlabel={\(x\)},
    ylabel={\(y\)},
    xtick={0,0.5,1,1.5,2},
    ytick={0,1,2,3,4},
    grid=both,
    grid style={line width=.1pt, draw=gray!25},
]
% left rectangles for x^2 on [0,2], n=4
\draw[fill=gray!15, draw=gray!60] (axis cs:0,0) rectangle (axis cs:0.5,0);
\draw[fill=gray!15, draw=gray!60] (axis cs:0.5,0) rectangle (axis cs:1,0.25);
\draw[fill=gray!15, draw=gray!60] (axis cs:1,0) rectangle (axis cs:1.5,1);
\draw[fill=gray!15, draw=gray!60] (axis cs:1.5,0) rectangle (axis cs:2,2.25);

\addplot[very thick, domain=0:2, samples=100] {x^2};
\node[anchor=west] at (axis cs:1.15,3.45) {\(f(x)=x^2\)};
\node[anchor=west] at (axis cs:0.15,0.65) {left rectangles};
\end{axis}
\end{tikzpicture}
\caption{Left endpoint rectangles estimate the area under \(y=x^2\) on \([0,2]\).}
\label{fig:left-rectangles-x-squared}
\end{figure}

Because \(x^2\) is increasing on \([0,2]\), every left endpoint rectangle lies below the
curve on its interval.  So this estimate is too small.

If we use right endpoints, the rectangles lie above the curve.  That estimate will be
too large.

This gives a useful strategy:

\[
\text{trap the unknown area between a lower estimate and an upper estimate.}
\]

\subsection{Left, right, and midpoint sums}
\label{subsec:left-right-midpoint-sums}

The choice of rectangle height matters.  Three common choices are left endpoints,
right endpoints, and midpoints.

For \(f(x)=x^2\) on \([0,2]\), with four equal subintervals and \(\Delta x=0.5\), the
left estimate is

\[
L_4
=
0.5\left(f(0)+f(0.5)+f(1)+f(1.5)\right).
\]

We computed

\[
L_4=1.75.
\]

The right estimate is

\[
R_4
=
0.5\left(f(0.5)+f(1)+f(1.5)+f(2)\right).
\]

So

\[
R_4
=
0.5\left(0.5^2+1^2+1.5^2+2^2\right)
=
0.5(0.25+1+2.25+4).
\]

Thus

\[
R_4=3.75.
\]

The midpoint estimate uses the midpoint of each subinterval:

\[
0.25,\qquad 0.75,\qquad 1.25,\qquad 1.75.
\]

So

\[
M_4
=
0.5\left(f(0.25)+f(0.75)+f(1.25)+f(1.75)\right).
\]

For \(f(x)=x^2\),

\[
M_4
=
0.5\left(0.25^2+0.75^2+1.25^2+1.75^2\right).
\]

Compute:

\[
M_4
=
0.5(0.0625+0.5625+1.5625+3.0625)
=
0.5(5.25)
=
2.625.
\]

So the three estimates are

\[
\boxed{
L_4=1.75,\qquad M_4=2.625,\qquad R_4=3.75.
}
\]

For an increasing function, the left estimate is below the true area and the right
estimate is above it.  The midpoint estimate often does better, but it is not guaranteed
to be an upper or lower estimate.

\begin{quote}
\textbf{Rectangular sums.}
A rectangular estimate has the form

\[
\text{area estimate}
=
(\text{height}_1)(\text{width}_1)
+
(\text{height}_2)(\text{width}_2)
+
\cdots.
\]

The height is a sampled value of the function.  The width is the length of the
corresponding subinterval.
\end{quote}

The word ``sum'' is important.  The integral will be a limit of these sums.

\subsection{Upper and lower sums}
\label{subsec:upper-lower-sums}

Left and right rectangles are easy to compute, but they are not always upper or lower
estimates.  If a function decreases, the left rectangles lie above the curve and the right
rectangles lie below it.  If a function rises and falls, both can be mixed.

A more reliable idea is to use the largest and smallest function values on each small
interval.

On a subinterval, define:

\[
\text{lower height}=\text{smallest value of }f\text{ on that subinterval},
\]

\[
\text{upper height}=\text{largest value of }f\text{ on that subinterval}.
\]

Then:

\[
\text{lower rectangle area}
=
(\text{lower height})(\text{width}),
\]

and

\[
\text{upper rectangle area}
=
(\text{upper height})(\text{width}).
\]

Adding all the lower rectangle areas gives a lower sum.  Adding all the upper
rectangle areas gives an upper sum.

For a continuous function, the Extreme Value Theorem guarantees that each small
closed interval really has a largest and smallest value.  This is one reason the
theorems about continuity matter.

\begin{figure}[h]
\centering
\begin{minipage}{0.48\textwidth}
\centering
\begin{tikzpicture}
\begin{axis}[
    width=\textwidth,
    height=0.72\textwidth,
    xmin=0, xmax=2.15,
    ymin=0, ymax=4.3,
    axis lines=left,
    xlabel={\(x\)},
    ylabel={},
    xtick={0,1,2},
    ytick={0,2,4},
    grid=both,
    grid style={line width=.1pt, draw=gray!25},
    title={lower sum}
]
\draw[fill=gray!15, draw=gray!60] (axis cs:0,0) rectangle (axis cs:0.5,0);
\draw[fill=gray!15, draw=gray!60] (axis cs:0.5,0) rectangle (axis cs:1,0.25);
\draw[fill=gray!15, draw=gray!60] (axis cs:1,0) rectangle (axis cs:1.5,1);
\draw[fill=gray!15, draw=gray!60] (axis cs:1.5,0) rectangle (axis cs:2,2.25);
\addplot[very thick, domain=0:2, samples=100] {x^2};
\end{axis}
\end{tikzpicture}
\end{minipage}
\hfill
\begin{minipage}{0.48\textwidth}
\centering
\begin{tikzpicture}
\begin{axis}[
    width=\textwidth,
    height=0.72\textwidth,
    xmin=0, xmax=2.15,
    ymin=0, ymax=4.3,
    axis lines=left,
    xlabel={\(x\)},
    ylabel={},
    xtick={0,1,2},
    ytick={0,2,4},
    grid=both,
    grid style={line width=.1pt, draw=gray!25},
    title={upper sum}
]
\draw[fill=gray!25, draw=gray!60] (axis cs:0,0) rectangle (axis cs:0.5,0.25);
\draw[fill=gray!25, draw=gray!60] (axis cs:0.5,0) rectangle (axis cs:1,1);
\draw[fill=gray!25, draw=gray!60] (axis cs:1,0) rectangle (axis cs:1.5,2.25);
\draw[fill=gray!25, draw=gray!60] (axis cs:1.5,0) rectangle (axis cs:2,4);
\addplot[very thick, domain=0:2, samples=100] {x^2};
\end{axis}
\end{tikzpicture}
\end{minipage}
\caption{For an increasing function, left rectangles give a lower sum and right rectangles give an upper sum.}
\label{fig:lower-upper-x-squared}
\end{figure}

For \(f(x)=x^2\) on \([0,2]\), the function is increasing.  So the lower sum with four
equal rectangles is the left estimate:

\[
1.75.
\]

The upper sum with four equal rectangles is the right estimate:

\[
3.75.
\]

Thus the true area \(A\), whatever it is, must satisfy

\[
1.75\le A\le 3.75.
\]

This is not very sharp.  But if we use more rectangles, the trap tightens.

\[
\begin{array}{c|c|c|c}
\text{number of rectangles} & \text{lower sum} & \text{upper sum} & \text{gap} \\ \hline
4 & 1.75 & 3.75 & 2 \\
8 & 2.1875 & 3.1875 & 1 \\
16 & 2.421875 & 2.921875 & 0.5
\end{array}
\]

The lower sums rise.  The upper sums fall.  The area is squeezed between them.

For \(f(x)=x^2\) on \([0,2]\), if we use \(n\) equal intervals, the width is

\[
\Delta x=\frac{2}{n}.
\]

Because the function is increasing, the upper sum exceeds the lower sum by

\[
\Delta x\bigl(f(2)-f(0)\bigr).
\]

Here that is

\[
\frac{2}{n}(4-0)=\frac{8}{n}.
\]

As \(n\) grows, the gap

\[
\frac{8}{n}
\]

approaches \(0\).  Therefore the lower and upper estimates are being forced toward
one common number.  That common number is the area.

\begin{quote}
\textbf{The squeezing idea for area.}
If upper sums and lower sums can be made as close as we want, then there is only one
reasonable number for the area between them.
\end{quote}

This is the geometric ancestor of the definite integral.

\subsection{Area as a limiting process}
\label{subsec:area-as-limiting-process}

A curved area is not found by one rectangle.  It is found by a limiting process.

\[
\text{area}
=
\lim_{\text{rectangle widths}\to0}
\text{rectangular sums}.
\]

This sentence is not yet a formal definition.  The formal definition will come after we
develop better notation.  But the idea is already clear:

\[
\text{thin rectangles make the curved region measurable.}
\]

It is helpful to test this idea on a region whose area we already know.

Consider

\[
f(x)=x
\]

on the interval \([0,2]\).  The region under the graph is a right triangle with base
\(2\) and height \(2\).  Its area is

\[
\frac12(2)(2)=2.
\]

Now see how rectangles find the same number.

Divide \([0,2]\) into \(n\) equal intervals.  Then

\[
\Delta x=\frac{2}{n}.
\]

The left endpoints are

\[
0,\quad \frac{2}{n},\quad \frac{4}{n},\quad \ldots,\quad \frac{2(n-1)}{n}.
\]

The left rectangle estimate is

\[
L_n
=
\frac{2}{n}
\left[
0+\frac{2}{n}+\frac{4}{n}+\cdots+\frac{2(n-1)}{n}
\right].
\]

Factor out \(\frac{2}{n}\) from the bracket:

\[
L_n
=
\frac{2}{n}\cdot \frac{2}{n}
\left[
0+1+2+\cdots+(n-1)
\right].
\]

Using

\[
0+1+2+\cdots+(n-1)=\frac{n(n-1)}{2},
\]

we get

\[
L_n
=
\frac{4}{n^2}\cdot \frac{n(n-1)}{2}
=
2\cdot\frac{n-1}{n}.
\]

As \(n\to\infty\),

\[
2\cdot\frac{n-1}{n}\to2.
\]

The right endpoint estimate is

\[
R_n
=
\frac{2}{n}
\left[
\frac{2}{n}+\frac{4}{n}+\cdots+\frac{2n}{n}
\right].
\]

So

\[
R_n
=
\frac{2}{n}\cdot \frac{2}{n}
\left[
1+2+\cdots+n
\right].
\]

Using

\[
1+2+\cdots+n=\frac{n(n+1)}{2},
\]

we get

\[
R_n
=
\frac{4}{n^2}\cdot \frac{n(n+1)}{2}
=
2\cdot\frac{n+1}{n}.
\]

As \(n\to\infty\),

\[
2\cdot\frac{n+1}{n}\to2.
\]

Thus

\[
L_n\to2
\qquad\text{and}\qquad
R_n\to2.
\]

The rectangle sums recover the triangle area.

\begin{figure}[h]
\centering
\begin{tikzpicture}
\begin{axis}[
    width=0.78\textwidth,
    height=0.48\textwidth,
    xmin=0, xmax=2.2,
    ymin=0, ymax=2.3,
    axis lines=left,
    xlabel={\(x\)},
    ylabel={\(y\)},
    xtick={0,1,2},
    ytick={0,1,2},
    grid=both,
    grid style={line width=.1pt, draw=gray!25},
]
% 8 left rectangles for f=x on [0,2]
\pgfplotsinvokeforeach{0,...,7}{
    \draw[fill=gray!15, draw=gray!55]
      (axis cs: {0.25*#1}, 0) rectangle (axis cs: {0.25*(#1+1)}, {0.25*#1});
}
\addplot[very thick, domain=0:2, samples=2] {x};
\node[anchor=west] at (axis cs:1.05,1.75) {\(y=x\)};
\node[anchor=west] at (axis cs:0.1,2.1) {rectangles approach the triangle area};
\end{axis}
\end{tikzpicture}
\caption{For \(y=x\) on \([0,2]\), rectangular sums approach the known triangle area \(2\).}
\label{fig:rectangles-triangle-area}
\end{figure}

For \(y=x\), geometry already knew the answer.  Rectangles confirmed it.

For \(y=x^2\), geometry does not give such an immediate answer.  Rectangles become
the method itself.  Later, the Fundamental Theorem of Calculus will give a fast way
to compute the same area exactly.

For now, the main point is this:

\[
\boxed{
\text{Area under a graph is an accumulation of thin rectangular pieces.}
}
\]

There is one caution.  Not every wild function behaves well enough for this process.
If a function jumps too violently at every scale, upper and lower sums may refuse to
come together.  Continuous functions do behave well, and we will return to that fact
later.

The next obstacle is notation.  Rectangular sums quickly become long:

\[
\frac{2}{n}
\left[
0+\frac{2}{n}+\frac{4}{n}+\cdots+\frac{2(n-1)}{n}
\right].
\]

To define the integral cleanly, we need a compact language for finite sums.  That is
where sigma notation enters.
% Requires in the preamble:
% \usepackage{amsmath,amssymb}
% \usepackage{tikz,pgfplots}
% \pgfplotsset{compat=1.18}

\section{Sigma notation and finite sums}
\label{sec:sigma-notation-finite-sums}

The last section ended with sums that quickly became too long to write by hand.  Even
a simple rectangular estimate can look like

\[
\frac{2}{n}
\left[
0+\frac{2}{n}+\frac{4}{n}+\cdots+\frac{2(n-1)}{n}
\right].
\]

The pattern is clear, but the notation is heavy.

Calculus needs a compact way to say:

\[
\text{add all the terms of this pattern.}
\]

That compact language is \emph{sigma notation}.  It is not a new idea.  It is just a
short name for a finite sum.

\subsection{Summation notation}
\label{subsec:summation-notation}

Suppose we want to write

\[
1^2+2^2+3^2+4^2+5^2.
\]

Instead of writing every term, we write

\[
\sum_{i=1}^{5} i^2.
\]

The symbol \(\sum\) is the Greek capital letter sigma.  It stands for ``sum.''

The notation

\[
\sum_{i=1}^{5} i^2
\]

means:

\[
\text{let \(i\) take the values }1,2,3,4,5,\text{ and add the results.}
\]

So

\[
\sum_{i=1}^{5} i^2
=
1^2+2^2+3^2+4^2+5^2
=
1+4+9+16+25
=
55.
\]

\begin{quote}
\textbf{Sigma notation.}
If \(m\le n\) are integers, then

\[
\boxed{
\sum_{i=m}^{n} a_i
=
a_m+a_{m+1}+a_{m+2}+\cdots+a_n.
}
\]

The letter \(i\) is called the \emph{index of summation}.  The number \(m\) is the
lower limit, and \(n\) is the upper limit.
\end{quote}

\paragraph{Example.}
Expand and evaluate

\[
\sum_{k=1}^{4}(2k-1).
\]

As \(k\) runs from \(1\) to \(4\), the terms are

\[
2(1)-1,\qquad 2(2)-1,\qquad 2(3)-1,\qquad 2(4)-1.
\]

So

\[
\sum_{k=1}^{4}(2k-1)
=
1+3+5+7
=
16.
\]

\paragraph{Example.}
Expand

\[
\sum_{j=0}^{3} 2^j.
\]

The index begins at \(0\), so

\[
\sum_{j=0}^{3} 2^j
=
2^0+2^1+2^2+2^3
=
1+2+4+8
=
15.
\]

\paragraph{The index letter is a placeholder.}
The sums

\[
\sum_{i=1}^{5} i^2,
\qquad
\sum_{k=1}^{5} k^2,
\qquad
\sum_{r=1}^{5} r^2
\]

all mean the same thing:

\[
1^2+2^2+3^2+4^2+5^2.
\]

The index letter is temporary.  It disappears after the sum is evaluated.  This is why it
is sometimes called a \emph{dummy index}.

This is similar to the variable in a function rule.  The expressions

\[
f(x)=x^2
\qquad\text{and}\qquad
f(t)=t^2
\]

describe the same function.  The letter changes, but the rule does not.

\paragraph{A warning about endpoints.}
The number of terms in

\[
\sum_{i=m}^{n} a_i
\]

is not \(n-m\).  It is

\[
n-m+1.
\]

For example,

\[
\sum_{i=3}^{7} a_i
=
a_3+a_4+a_5+a_6+a_7
\]

has \(5\) terms, and

\[
7-3+1=5.
\]

This small detail matters when we count rectangles.

\subsection{Sigma notation for rectangular sums}
\label{subsec:sigma-notation-rectangular-sums}

Return to the left rectangle estimate for

\[
f(x)=x^2
\]

on \([0,2]\) with four equal subintervals.  The width is

\[
\Delta x=\frac12.
\]

The left endpoints are

\[
0,\qquad \frac12,\qquad 1,\qquad \frac32.
\]

We previously wrote the estimate as

\[
\frac12\left[
0^2+\left(\frac12\right)^2+1^2+\left(\frac32\right)^2
\right].
\]

Using sigma notation, the same sum is

\[
\sum_{i=0}^{3}
\left(\frac{i}{2}\right)^2
\left(\frac12\right).
\]

Here

\[
\frac{i}{2}
\]

runs through

\[
0,\qquad \frac12,\qquad 1,\qquad \frac32
\]

as \(i\) runs from \(0\) to \(3\).

This may not look much shorter for four rectangles.  But for \(n\) rectangles, the
advantage is enormous.

Divide \([0,2]\) into \(n\) equal subintervals.  Then

\[
\Delta x=\frac{2}{n}.
\]

The right endpoints are

\[
\frac{2}{n},\quad \frac{4}{n},\quad \frac{6}{n},\quad \ldots,\quad \frac{2n}{n}.
\]

The right rectangle sum for \(f(x)=x^2\) is

\[
\sum_{i=1}^{n}
\left(\frac{2i}{n}\right)^2
\left(\frac{2}{n}\right).
\]

This one line represents all \(n\) rectangles.

\[
\boxed{
\text{Sigma notation lets us write a finite rectangular sum without writing every rectangle.}
}
\]

\subsection{Arithmetic sums}
\label{subsec:arithmetic-sums}

Some finite sums have simple formulas.  The first one is the sum of the first \(n\)
positive integers:

\[
1+2+3+\cdots+n.
\]

Write

\[
S=1+2+3+\cdots+n.
\]

Now write the same sum backward:

\[
S=n+(n-1)+(n-2)+\cdots+1.
\]

Add the two lines term by term:

\[
\begin{aligned}
2S
&=
(1+n)+(2+n-1)+(3+n-2)+\cdots+(n+1)\\
&=
(n+1)+(n+1)+(n+1)+\cdots+(n+1).
\end{aligned}
\]

There are \(n\) copies of \(n+1\), so

\[
2S=n(n+1).
\]

Therefore

\[
\boxed{
\sum_{i=1}^{n} i
=
1+2+3+\cdots+n
=
\frac{n(n+1)}{2}.
}
\]

\paragraph{Example.}
Find

\[
1+2+3+\cdots+100.
\]

Using the formula,

\[
\sum_{i=1}^{100} i
=
\frac{100(101)}{2}
=
5050.
\]

\paragraph{Example.}
Evaluate

\[
\sum_{i=1}^{n}(3i+2).
\]

Use the fact that sums distribute over addition:

\[
\sum_{i=1}^{n}(3i+2)
=
\sum_{i=1}^{n}3i+\sum_{i=1}^{n}2.
\]

Now pull out the constant \(3\):

\[
=
3\sum_{i=1}^{n}i+\sum_{i=1}^{n}2.
\]

Since there are \(n\) copies of \(2\),

\[
\sum_{i=1}^{n}2=2n.
\]

Therefore

\[
\sum_{i=1}^{n}(3i+2)
=
3\cdot \frac{n(n+1)}{2}+2n.
\]

So

\[
\boxed{
\sum_{i=1}^{n}(3i+2)
=
\frac{3n(n+1)}{2}+2n.
}
\]

We can simplify further if we want:

\[
\frac{3n(n+1)}{2}+2n
=
\frac{3n^2+3n+4n}{2}
=
\frac{3n^2+7n}{2}.
\]

\subsection{Basic rules for finite sums}
\label{subsec:basic-rules-for-finite-sums}

The example used rules that are worth naming.

\begin{quote}
\textbf{Rules for finite sums.}
For constants \(c\) and finite lists \(a_i,b_i\),

\[
\boxed{
\sum_{i=m}^{n} c
=
c(n-m+1),
}
\]

\[
\boxed{
\sum_{i=m}^{n} c a_i
=
c\sum_{i=m}^{n} a_i,
}
\]

and

\[
\boxed{
\sum_{i=m}^{n} (a_i+b_i)
=
\sum_{i=m}^{n} a_i+\sum_{i=m}^{n} b_i.
}
\]
\end{quote}

These rules are not mysterious.  They come from ordinary addition.

For example,

\[
\sum_{i=1}^{n}(a_i+b_i)
\]

means

\[
(a_1+b_1)+(a_2+b_2)+\cdots+(a_n+b_n).
\]

Regroup the terms:

\[
(a_1+a_2+\cdots+a_n)+(b_1+b_2+\cdots+b_n).
\]

That is

\[
\sum_{i=1}^{n} a_i+\sum_{i=1}^{n} b_i.
\]

A finite sum behaves like a long addition problem because that is exactly what it is.

\subsection{Geometric sums}
\label{subsec:geometric-sums}

Another important pattern is

\[
1+r+r^2+r^3+\cdots+r^{n-1}.
\]

This is called a \emph{geometric sum}.  Each term is obtained from the previous term
by multiplying by \(r\).

Let

\[
G=1+r+r^2+\cdots+r^{n-1}.
\]

Multiply by \(r\):

\[
rG=r+r^2+r^3+\cdots+r^n.
\]

Subtract:

\[
G-rG
=
(1+r+r^2+\cdots+r^{n-1})
-
(r+r^2+\cdots+r^{n-1}+r^n).
\]

Most terms cancel, leaving

\[
G-rG=1-r^n.
\]

So

\[
(1-r)G=1-r^n.
\]

If \(r\neq1\), then

\[
\boxed{
1+r+r^2+\cdots+r^{n-1}
=
\frac{1-r^n}{1-r}.
}
\]

In sigma notation,

\[
\boxed{
\sum_{i=0}^{n-1} r^i
=
\frac{1-r^n}{1-r},
\qquad r\neq1.
}
\]

If \(r=1\), then every term is \(1\), so

\[
\sum_{i=0}^{n-1}1^i=n.
\]

\paragraph{Example.}
Evaluate

\[
1+\frac12+\frac14+\frac18+\cdots+\frac{1}{2^{n-1}}.
\]

Here

\[
r=\frac12.
\]

So

\[
\sum_{i=0}^{n-1}\left(\frac12\right)^i
=
\frac{1-\left(\frac12\right)^n}{1-\frac12}.
\]

Since

\[
1-\frac12=\frac12,
\]

we get

\[
\sum_{i=0}^{n-1}\left(\frac12\right)^i
=
2\left(1-\frac{1}{2^n}\right).
\]

As \(n\) grows, the term \(\frac{1}{2^n}\) becomes very small.  So the sum approaches

\[
2.
\]

This is a preview of infinite series.  A finite geometric sum has an exact formula.  When
the number of terms grows without bound, the finite formula tells us what the infinite
process approaches.

\subsection{Sums of powers}
\label{subsec:sums-of-powers}

Rectangular sums often contain powers of the index:

\[
\sum_{i=1}^{n} i,
\qquad
\sum_{i=1}^{n} i^2,
\qquad
\sum_{i=1}^{n} i^3.
\]

The first formula is already known:

\[
\sum_{i=1}^{n} i=\frac{n(n+1)}{2}.
\]

The next two are:

\[
\boxed{
\sum_{i=1}^{n} i^2
=
\frac{n(n+1)(2n+1)}{6},
}
\]

and

\[
\boxed{
\sum_{i=1}^{n} i^3
=
\left(\frac{n(n+1)}{2}\right)^2.
}
\]

Together:

\[
\begin{array}{c|c}
\textbf{sum} & \textbf{formula} \\ \hline
\displaystyle \sum_{i=1}^{n} 1 & n \\[8pt]
\displaystyle \sum_{i=1}^{n} i & \displaystyle \frac{n(n+1)}{2} \\[10pt]
\displaystyle \sum_{i=1}^{n} i^2 & \displaystyle \frac{n(n+1)(2n+1)}{6} \\[10pt]
\displaystyle \sum_{i=1}^{n} i^3 & \displaystyle \left(\frac{n(n+1)}{2}\right)^2
\end{array}
\]

These formulas are useful because powers of \(i\) appear when we evaluate Riemann
sums for functions such as \(x\), \(x^2\), and \(x^3\).

\paragraph{Why the sum of squares formula is the right one.}
One way to find the formula for

\[
\sum_{i=1}^{n} i^2
\]

is to use a telescoping sum.

Start with the identity

\[
(k+1)^3-k^3=3k^2+3k+1.
\]

Now add this identity from \(k=1\) to \(k=n\):

\[
\sum_{k=1}^{n}\bigl((k+1)^3-k^3\bigr)
=
\sum_{k=1}^{n}(3k^2+3k+1).
\]

The left side telescopes:

\[
(2^3-1^3)+(3^3-2^3)+\cdots+((n+1)^3-n^3)
=
(n+1)^3-1.
\]

The right side is

\[
3\sum_{k=1}^{n}k^2
+
3\sum_{k=1}^{n}k
+
\sum_{k=1}^{n}1.
\]

So

\[
(n+1)^3-1
=
3\sum_{k=1}^{n}k^2
+
3\cdot \frac{n(n+1)}{2}
+
n.
\]

Solving for \(\sum k^2\) gives

\[
\sum_{k=1}^{n}k^2
=
\frac{n(n+1)(2n+1)}{6}.
\]

The formula comes from adding finite differences.

This is another hint of the larger story:

\[
\text{differences telescope, and sums accumulate.}
\]

\paragraph{Example.}
Evaluate

\[
\sum_{i=1}^{n}(4i^2-3i+2).
\]

Use the sum rules:

\[
\sum_{i=1}^{n}(4i^2-3i+2)
=
4\sum_{i=1}^{n}i^2
-
3\sum_{i=1}^{n}i
+
2\sum_{i=1}^{n}1.
\]

Substitute the formulas:

\[
=
4\cdot \frac{n(n+1)(2n+1)}{6}
-
3\cdot \frac{n(n+1)}{2}
+
2n.
\]

So

\[
\boxed{
\sum_{i=1}^{n}(4i^2-3i+2)
=
\frac{2n(n+1)(2n+1)}{3}
-
\frac{3n(n+1)}{2}
+
2n.
}
\]

We could simplify the expression, but in calculus it is often more useful to see which
standard sums produced it.

\paragraph{Iterating for higher powers.}
This telescoping technique generalizes to any integer power \(p \ge 1\). To find a closed-form formula for \(\sum_{i=1}^{n} i^p\), we start with the binomial expansion of the difference of \((p+1)\)-th powers:

\[
(k+1)^{p+1} - k^{p+1}
=
(p+1)k^p + \binom{p+1}{2}k^{p-1} + \cdots + \binom{p+1}{p}k + 1.
\]

Summing both sides from \(k=1\) to \(n\) telescopes the left side entirely:

\[
(n+1)^{p+1}-1
=
(p+1)\sum_{k=1}^{n}k^p 
+ 
\binom{p+1}{2}\sum_{k=1}^{n}k^{p-1} 
+ \cdots + 
\sum_{k=1}^{n}1.
\]

Because the right side isolates \((p+1)\sum k^p\) alongside a linear combination of sums of strictly lower degree, we can solve for \(\sum k^p\) as long as we have already computed the formulas for all lower powers.
\subsection{From finite sums to integrals}
\label{subsec:finite-sums-to-integrals}

Sigma notation was introduced because rectangular sums need it.

Divide the interval \([a,b]\) into \(n\) equal subintervals.  Then the common width is

\[
\Delta x=\frac{b-a}{n}.
\]

If we use right endpoints, then the \(i\)th sample point is

\[
x_i=a+i\Delta x.
\]

The right rectangle sum is

\[
\sum_{i=1}^{n} f(x_i)\Delta x
=
\sum_{i=1}^{n} f(a+i\Delta x)\Delta x.
\]

This is the basic finite object that will become an integral.

\paragraph{Example: \(y=x\) on \([0,2]\).}
Here

\[
f(x)=x,
\qquad
a=0,
\qquad
b=2,
\qquad
\Delta x=\frac{2}{n}.
\]

The right endpoints are

\[
x_i=\frac{2i}{n}.
\]

So the right rectangle sum is

\[
R_n
=
\sum_{i=1}^{n}
\left(\frac{2i}{n}\right)
\left(\frac{2}{n}\right).
\]

Pull out the constants:

\[
R_n
=
\frac{4}{n^2}\sum_{i=1}^{n} i.
\]

Use

\[
\sum_{i=1}^{n} i=\frac{n(n+1)}{2}.
\]

Then

\[
R_n
=
\frac{4}{n^2}\cdot \frac{n(n+1)}{2}
=
2\cdot \frac{n+1}{n}.
\]

As \(n\to\infty\),

\[
2\cdot \frac{n+1}{n}\to2.
\]

So the right rectangle sums approach

\[
2.
\]

That matches the triangle area under \(y=x\) from \(0\) to \(2\).

\paragraph{Example: \(y=x^2\) on \([0,2]\).}
This is the curved region from the last section.

Again

\[
\Delta x=\frac{2}{n},
\qquad
x_i=\frac{2i}{n}.
\]

The right rectangle sum is

\[
R_n
=
\sum_{i=1}^{n}
\left(\frac{2i}{n}\right)^2
\left(\frac{2}{n}\right).
\]

Simplify:

\[
R_n
=
\sum_{i=1}^{n}
\frac{4i^2}{n^2}\cdot \frac{2}{n}
=
\frac{8}{n^3}\sum_{i=1}^{n}i^2.
\]

Use the sum of squares formula:

\[
\sum_{i=1}^{n}i^2
=
\frac{n(n+1)(2n+1)}{6}.
\]

Then

\[
R_n
=
\frac{8}{n^3}\cdot
\frac{n(n+1)(2n+1)}{6}.
\]

So

\[
R_n
=
\frac{4}{3}\cdot
\frac{(n+1)(2n+1)}{n^2}.
\]

Expand the numerator:

\[
(n+1)(2n+1)=2n^2+3n+1.
\]

Therefore

\[
R_n
=
\frac{4}{3}
\left(
2+\frac{3}{n}+\frac{1}{n^2}
\right).
\]

As \(n\to\infty\),

\[
R_n\to
\frac{4}{3}\cdot2
=
\frac{8}{3}.
\]

The exact area under \(y=x^2\) from \(0\) to \(2\) should therefore be

\[
\frac{8}{3}.
\]

The next section gives this limiting number a name and a notation.

\[
\boxed{
\text{The definite integral is the limit of Riemann sums.}
}
\]

\section{The definite integral}
\label{sec:definite-integral}

Sigma notation made rectangular sums readable.  Now we give the limiting process a
name.

The definite integral is the number that rectangular sums approach when every
rectangle becomes thin.

\[
\text{finite sum of rectangles}
\quad\longrightarrow\quad
\text{definite integral}.
\]

This definition is the formal version of the accumulation idea that began the chapter.

\subsection{Partitions and sample points}
\label{subsec:partitions-sample-points}

To build rectangular sums carefully, we first divide an interval into pieces.

Let

\[
[a,b]
\]

be a closed interval.  A \emph{partition} \(P\) of \([a,b]\) is a finite list of points

\[
a=x_0<x_1<x_2<\cdots<x_n=b.
\]

The subintervals are

\[
[x_0,x_1],
\quad
[x_1,x_2],
\quad
\ldots,
\quad
[x_{n-1},x_n].
\]

The width of the \(i\)th subinterval is

\[
\Delta x_i=x_i-x_{i-1}.
\]

If all widths are equal, then

\[
\Delta x_i=\Delta x=\frac{b-a}{n}.
\]

But equal widths are not required.

The largest width is called the \emph{mesh} of the partition:

\[
\lVert P\rVert
=
\max\{\Delta x_1,\Delta x_2,\ldots,\Delta x_n\}.
\]

The condition

\[
\lVert P\rVert\to0
\]

means that every subinterval is becoming small.

On each subinterval \([x_{i-1},x_i]\), choose a sample point

\[
x_i^*
\]

inside that subinterval:

\[
x_i^*\in[x_{i-1},x_i].
\]

The sample point determines the height of the rectangle:

\[
f(x_i^*).
\]

\begin{figure}[h]
\centering
\begin{tikzpicture}
\begin{axis}[
    width=0.82\textwidth,
    height=0.48\textwidth,
    xmin=0, xmax=3.15,
    ymin=0, ymax=4.3,
    axis lines=left,
    xlabel={\(x\)},
    ylabel={\(y\)},
    xtick={0,0.6,1.25,2.1,3},
    xticklabels={\(x_0\),\(x_1\),\(x_2\),\(x_3\),\(x_4\)},
    ytick={0,1,2,3,4},
    grid=both,
    grid style={line width=.1pt, draw=gray!25},
]
% Function f(x)=1+0.25x+0.35x^2
\foreach \a/\b/\s in {0/0.6/0.35, 0.6/1.25/0.95, 1.25/2.1/1.65, 2.1/3/2.55} {
    \pgfmathsetmacro{\height}{1+0.25*\s+0.35*\s*\s}
    \edef\temp{
        \noexpand\draw[fill=gray!15, draw=gray!60]
            (axis cs:\a,0) rectangle (axis cs:\b,\height);
        \noexpand\addplot[only marks, mark=*] coordinates {(\s,\height)};
    }
    \temp
}
\addplot[very thick, domain=0:3, samples=150] {1+0.25*x+0.35*x^2};
\node[anchor=west] at (axis cs:1.7,3.65) {\(y=f(x)\)};
\node[anchor=west] at (axis cs:0.15,0.45) {sample-point rectangles};
\end{axis}
\end{tikzpicture}
\caption{A partition divides the interval. A sample point in each subinterval determines the rectangle height.}
\label{fig:partition-sample-points}
\end{figure}

The sample point may be the left endpoint:

\[
x_i^*=x_{i-1}.
\]

It may be the right endpoint:

\[
x_i^*=x_i.
\]

It may be the midpoint:

\[
x_i^*=\frac{x_{i-1}+x_i}{2}.
\]

Or it may be some other point in the subinterval.  The definition of the integral will
ask whether all reasonable choices are forced toward the same number as the partition
becomes fine.

\subsection{Riemann sums}
\label{subsec:riemann-sums}

For each subinterval, the rectangle has width

\[
\Delta x_i
\]

and height

\[
f(x_i^*).
\]

So its signed rectangular contribution is

\[
f(x_i^*)\Delta x_i.
\]

Adding all the contributions gives a \emph{Riemann sum}.

\begin{quote}
\textbf{Riemann sum.}
Given a partition

\[
P:\quad a=x_0<x_1<\cdots<x_n=b
\]

and sample points

\[
x_i^*\in[x_{i-1},x_i],
\]

the corresponding Riemann sum for \(f\) on \([a,b]\) is

\[
\boxed{
\sum_{i=1}^{n} f(x_i^*)\Delta x_i.
}
\]
\end{quote}

If \(f(x)\ge0\), this sum approximates area under the graph.

If \(f\) can be negative, the sum still makes sense.  Terms below the axis contribute
negatively.  That signed interpretation will become central in the next section.

\paragraph{Example.}
Let

\[
f(x)=1+x^2
\]

on \([0,2]\).  Divide the interval into four equal subintervals and use right endpoints.

Then

\[
\Delta x=\frac{2-0}{4}=\frac12.
\]

The right endpoints are

\[
\frac12,\qquad 1,\qquad \frac32,\qquad 2.
\]

The Riemann sum is

\[
\sum_{i=1}^{4} f\left(\frac{i}{2}\right)\frac12.
\]

Compute:

\[
\begin{aligned}
\sum_{i=1}^{4} f\left(\frac{i}{2}\right)\frac12
&=
\frac12\left[
f\left(\frac12\right)+f(1)+f\left(\frac32\right)+f(2)
\right] \\[4pt]
&=
\frac12\left[
\left(1+\frac14\right)+(1+1)+\left(1+\frac94\right)+(1+4)
\right].
\end{aligned}
\]

Thus

\[
\begin{aligned}
\sum_{i=1}^{4} f\left(\frac{i}{2}\right)\frac12
&=
\frac12\left[
\frac54+2+\frac{13}{4}+5
\right] \\[4pt]
&=
\frac12\left(\frac{42}{4}\right) \\[4pt]
&=
\frac{21}{4}.
\end{aligned}
\]

So this particular Riemann sum is

\[
\boxed{\frac{21}{4}=5.25.}
\]

It is one approximation.  The integral is what these approximations approach as the
subintervals become thin.

\subsection{Definition of the definite integral}
\label{subsec:definite-integral-definition}

We are ready for the definition.

\begin{quote}
\textbf{Definite integral.}
Let \(f\) be a bounded function on \([a,b]\).  If the Riemann sums

\[
\sum_{i=1}^{n} f(x_i^*)\Delta x_i
\]

approach one common number as

\[
\lVert P\rVert\to0,
\]

no matter how the partitions and sample points are chosen, then \(f\) is called
\emph{integrable} on \([a,b]\).

That common number is the \emph{definite integral} of \(f\) from \(a\) to \(b\), written

\[
\boxed{
\int_a^b f(x)\,dx
=
\lim_{\lVert P\rVert\to0}
\sum_{i=1}^{n} f(x_i^*)\Delta x_i.
}
\]
\end{quote}

The notation is full of memory.

\[
\int
\]

is an elongated \(S\), reminding us of a sum.

The expression

\[
f(x)
\]

reminds us of the rectangle height.

The symbol

\[
dx
\]

reminds us of the small horizontal width, the limiting descendant of \(\Delta x_i\).

So

\[
\int_a^b f(x)\,dx
\]

means:

\[
\text{accumulate \(f(x)\) along the interval from \(a\) to \(b\).}
\]

The \(dx\) is not decoration.  It tells us which variable is being accumulated and
points to the small step in the input.  Later, when we reverse direction along the
interval, that orientation will matter:

\[
\int_b^a f(x)\,dx=-\int_a^b f(x)\,dx.
\]

For now, when \(a<b\), the widths \(\Delta x_i\) are positive.

\paragraph{The variable of integration is temporary.}
The expressions

\[
\int_a^b f(x)\,dx,
\qquad
\int_a^b f(t)\,dt,
\qquad
\int_a^b f(u)\,du
\]

mean the same number.  The variable is a placeholder inside the integral, just as the
index in a finite sum is a placeholder:

\[
\sum_{i=1}^{n} i^2
=
\sum_{k=1}^{n} k^2.
\]

The bounds \(a\) and \(b\) are not temporary.  They determine the interval of
accumulation.

\subsection{Equal-width Riemann sums}
\label{subsec:equal-width-riemann-sums}

Many calculations use equal subintervals.  If \([a,b]\) is divided into \(n\) equal
pieces, then

\[
\Delta x=\frac{b-a}{n}.
\]

Using right endpoints,

\[
x_i=a+i\Delta x,
\qquad i=1,2,\ldots,n.
\]

The right Riemann sum is

\[
\sum_{i=1}^{n} f(a+i\Delta x)\Delta x.
\]

If \(f\) is integrable, then this particular sequence of sums approaches the integral:

\[
\boxed{
\int_a^b f(x)\,dx
=
\lim_{n\to\infty}
\sum_{i=1}^{n} f(a+i\Delta x)\Delta x,
\qquad
\Delta x=\frac{b-a}{n}.
}
\]

This formula is not the definition in its full strength, because it uses only equal
widths and right endpoints.  But for integrable functions, it is often the most convenient
way to compute from the definition.

\subsection{Computing an integral from the definition}
\label{subsec:computing-integral-definition}

The area under

\[
y=x^2
\]

from \(0\) to \(2\) was estimated in the previous section.  Now we compute its exact
integral from Riemann sums.

Let

\[
f(x)=x^2,
\qquad
0\le x\le2.
\]

Divide \([0,2]\) into \(n\) equal subintervals.  Then

\[
\Delta x=\frac{2}{n}.
\]

Use right endpoints:

\[
x_i=\frac{2i}{n}.
\]

The right Riemann sum is

\[
R_n
=
\sum_{i=1}^{n}
f(x_i)\Delta x.
\]

Since \(f(x)=x^2\),

\[
R_n
=
\sum_{i=1}^{n}
\left(\frac{2i}{n}\right)^2
\left(\frac{2}{n}\right).
\]

Simplify:

\[
R_n
=
\sum_{i=1}^{n}
\frac{4i^2}{n^2}\cdot\frac{2}{n}
=
\frac{8}{n^3}\sum_{i=1}^{n}i^2.
\]

Use

\[
\sum_{i=1}^{n}i^2
=
\frac{n(n+1)(2n+1)}{6}.
\]

Then

\[
R_n
=
\frac{8}{n^3}\cdot\frac{n(n+1)(2n+1)}{6}.
\]

Thus

\[
R_n
=
\frac{4}{3}
\cdot
\frac{(n+1)(2n+1)}{n^2}.
\]

Expand:

\[
R_n
=
\frac{4}{3}
\left(
2+\frac{3}{n}+\frac{1}{n^2}
\right).
\]

Therefore

\[
\lim_{n\to\infty}R_n
=
\frac{4}{3}\cdot2
=
\frac{8}{3}.
\]

So

\[
\boxed{
\int_0^2 x^2\,dx=\frac{8}{3}.
}
\]

Because \(x^2\ge0\) on \([0,2]\), this number is also the area under the curve.

\begin{figure}[h]
\centering
\begin{tikzpicture}
\begin{axis}[
    width=0.78\textwidth,
    height=0.48\textwidth,
    xmin=0, xmax=2.15,
    ymin=0, ymax=4.4,
    axis lines=left,
    xlabel={\(x\)},
    ylabel={\(y\)},
    xtick={0,1,2},
    ytick={0,1,2,3,4},
    grid=both,
    grid style={line width=.1pt, draw=gray!25},
]
% right rectangles, n=8
\pgfplotsinvokeforeach{1,...,8}{
    \draw[fill=gray!15, draw=gray!55]
        (axis cs: {0.25*(#1-1)}, 0) rectangle (axis cs: {0.25*#1}, {(0.25*#1)^2});
}
\addplot[very thick, domain=0:2, samples=150] {x^2};
\node[anchor=west] at (axis cs:0.9,3.55) {\(y=x^2\)};
\node[anchor=west] at (axis cs:0.15,1.2) {right Riemann rectangles};
\end{axis}
\end{tikzpicture}
\caption{The exact area under \(y=x^2\) on \([0,2]\) is the limit of these rectangular sums: \(\int_0^2 x^2\,dx=8/3\).}
\label{fig:riemann-x-squared-right}
\end{figure}

This calculation is the first exact curved area in the chapter.  It came from three
ingredients:

\[
\text{partition},
\qquad
\text{Riemann sum},
\qquad
\text{limit}.
\]

\subsection{Why the sample point should not matter}
\label{subsec:why-sample-point-should-not-matter}

The definition of the integral asks for more than the convergence of one special sum.
It asks that all sufficiently fine Riemann sums approach the same number.

Why such a strong requirement?

Because the integral should belong to the function and the interval, not to our particular
choice of sample points.

For a smooth graph, this is natural.  If a subinterval is very narrow, the values of \(f\)
inside it are close to one another.  A left endpoint, a right endpoint, and a midpoint give
nearly the same height.  The rectangle heights may differ a little, but the differences
become negligible as the widths shrink.

For \(f(x)=x^2\) on \([0,2]\), the left and right sums have the same limit.  The right
sum is

\[
R_n=
\sum_{i=1}^{n}
\left(\frac{2i}{n}\right)^2
\left(\frac{2}{n}\right)
\to
\frac83.
\]

The left sum is

\[
L_n=
\sum_{i=0}^{n-1}
\left(\frac{2i}{n}\right)^2
\left(\frac{2}{n}\right).
\]

Using the same sum-of-squares formula, one finds

\[
L_n\to\frac83.
\]

The upper and lower sums also squeeze toward the same number.  That is why the
integral exists and equals

\[
\frac83.
\]

\[
\boxed{
\text{When all fine rectangular sums are forced toward one number, that number is the integral.}
}
\]

\subsection{Integrability, informally first}
\label{subsec:integrability-informally-first}

The definition says:

\[
\text{all fine Riemann sums must approach one common value.}
\]

This is a strong condition, but most functions we meet in first-year calculus satisfy it.

\begin{quote}
\textbf{Working theorem.}
If \(f\) is continuous on \([a,b]\), then \(f\) is integrable on \([a,b]\).
\end{quote}

We will discuss this more carefully later.  The intuition is that a continuous function
cannot jump wildly inside tiny intervals.  When the subintervals are small, the function
values inside each subinterval are forced to stay close together.

Continuity is not required, though.  A function with a finite number of jump
discontinuities can still be integrable.

For example, the step function

\[
H(x)=
\begin{cases}
0, & 0\le x<1,\\
1, & 1\le x\le2
\end{cases}
\]

is integrable on \([0,2]\).  The jump at \(x=1\) affects only one point.  Riemann sums
still settle down to the total area:

\[
1.
\]

But not every bounded function is integrable.

\paragraph{Warning example: too much oscillation.}
Define a function \(D\) on \([0,1]\) by

\[
D(x)=
\begin{cases}
1, & x\text{ is rational},\\
0, & x\text{ is irrational}.
\end{cases}
\]

Every interval, no matter how small, contains rational numbers and irrational numbers.
So on every subinterval, we can choose a rational sample point and get height \(1\), or
choose an irrational sample point and get height \(0\).

If every sample point is rational, the Riemann sum is

\[
\sum_{i=1}^{n}1\cdot \Delta x_i
=
\sum_{i=1}^{n}\Delta x_i
=
1.
\]

If every sample point is irrational, the Riemann sum is

\[
\sum_{i=1}^{n}0\cdot \Delta x_i
=
0.
\]

These two outcomes do not approach one common number.  They remain separated by
\(1\), no matter how fine the partition is.

Therefore \(D\) is not Riemann integrable on \([0,1]\).

\begin{quote}
\textbf{Moral.}
The definite integral is not just any limit of convenient sums.  It exists only when
fine Riemann sums are forced toward the same value, regardless of the sample points.
\end{quote}

This warning is not meant to frighten us.  It tells us why the definition is honest.
The continuous functions, piecewise continuous functions, and ordinary models of
science and geometry are integrable.  But the definition knows how to reject functions
that oscillate too violently at every scale.

We now have the central object of integral calculus:

\[
\int_a^b f(x)\,dx.
\]

When \(f\ge0\), it is area under the graph.  But the original motivation came from
velocity, and velocity can be negative.  The next step is to understand the integral as
\emph{signed} accumulation: area above the axis counts positively, area below the axis
counts negatively.
\section{Signed area and net change}
\label{sec:signed-area-net-change}

The definite integral was built from Riemann sums:

\[
\sum_{j=1}^{n} f(x_j^*)\Delta x_j.
\]

Each term has two parts.  The width \(\Delta x_j\) is positive when we move from left to
right.  The height \(f(x_j^*)\) can be positive, zero, or negative.

That one fact changes the meaning of area.

If the graph lies above the \(x\)-axis, the rectangles contribute positively.  If the graph
lies below the \(x\)-axis, the rectangles contribute negatively.  So the definite integral is
not ordinary area first.  It is \emph{signed accumulation}.

\[
\boxed{
\int_a^b f(x)\,dx
=
\text{area above the axis}
-
\text{area below the axis}.
}
\]

Ordinary area is never negative.  A definite integral can be negative, because it is designed
to remember direction.

\subsection{Area above the axis}
\label{subsec:area-above-axis}

When

\[
f(x)\ge 0
\]

on \([a,b]\), all the small rectangles in a Riemann sum have nonnegative height.  In that
case,

\[
\int_a^b f(x)\,dx
\]

is the ordinary area under the graph.

For example, suppose \(f(x)=2\) on \([0,5]\).  Then the graph is a horizontal line, and the
area is a rectangle:

\[
\int_0^5 2\,dx=2\cdot5=10.
\]

The integral is positive because the graph is above the axis.

More generally, if \(f(x)\ge0\), then the integral accumulates positive amount.  In a
motion problem, this might mean forward displacement.  In a density problem, it might
mean mass.  In a power problem, it might mean energy used.

The same rectangle picture from the beginning of the book is still present:

\[
\text{height}\cdot\text{width}
=
\text{rate}\cdot\text{input step}.
\]

The integral adds infinitely many such pieces.

\subsection{Area below the axis}
\label{subsec:area-below-axis}

When

\[
f(x)\le 0
\]

on \([a,b]\), all the small rectangles have negative height.  The ordinary geometric area
between the graph and the axis is positive, but the integral is negative.

For example,

\[
f(x)=-3
\]

on \([0,4]\) gives

\[
\int_0^4 -3\,dx=-3\cdot4=-12.
\]

The ordinary area of the rectangle is \(12\).  The signed area is \(-12\).

This is not a defect.  It is the point.

If \(v(t)=-3\) meters/second for \(4\) seconds, then the object moves \(12\) meters in
the negative direction.  Its displacement is

\[
-12\text{ meters}.
\]

So the signed integral matches the signed physical change.

\subsection{A graph with both signs}
\label{subsec:graph-with-both-signs}

Let

\[
f(x)=x-1
\]

on the interval \([0,3]\).

The graph crosses the \(x\)-axis at \(x=1\).  It lies below the axis on \([0,1]\), and above
the axis on \([1,3]\).

\begin{figure}[h]
\centering
\begin{tikzpicture}
\begin{axis}[
    width=0.78\textwidth,
    height=0.48\textwidth,
    xmin=0, xmax=3.2,
    ymin=-1.3, ymax=2.3,
    axis lines=middle,
    xlabel={\(x\)},
    ylabel={\(y\)},
    xtick={0,1,3},
    ytick={-1,0,2},
    grid=both,
    grid style={line width=.1pt, draw=gray!25},
]
\draw[fill=gray!25, draw=none]
  (axis cs:0,0) -- (axis cs:1,0) -- (axis cs:0,-1) -- cycle;
\draw[fill=gray!15, draw=none]
  (axis cs:1,0) -- (axis cs:3,0) -- (axis cs:3,2) -- cycle;
\addplot[very thick, domain=0:3, samples=2] {x-1};
\node[anchor=north] at (axis cs:0.45,-0.12) {\(-\frac12\)};
\node[anchor=south] at (axis cs:2.2,0.45) {\(+2\)};
\node[anchor=west] at (axis cs:2.05,1.65) {\(y=x-1\)};
\end{axis}
\end{tikzpicture}
\caption{The integral counts area below the axis negatively and area above the axis positively.}
\label{fig:signed-area-linear}
\end{figure}

The triangle below the axis has base \(1\) and height \(1\), so its ordinary area is

\[
\frac12.
\]

Because it lies below the axis, it contributes

\[
-\frac12
\]

to the integral.

The triangle above the axis has base \(2\) and height \(2\), so its ordinary area is

\[
\frac12(2)(2)=2.
\]

Therefore

\[
\int_0^3 (x-1)\,dx
=
-\frac12+2
=
\frac32.
\]

The integral is positive, even though part of the graph is below the axis, because the
positive part is larger than the negative part.

\subsection{Cancellation and meaning}
\label{subsec:cancellation-meaning}

Signed area can cancel.

For example, consider

\[
f(x)=x
\]

on \([-1,1]\).  The graph lies below the axis on \([-1,0]\) and above the axis on
\([0,1]\).  The two triangles have the same ordinary area:

\[
\frac12.
\]

The signed integral is

\[
\int_{-1}^{1} x\,dx
=
-\frac12+\frac12
=
0.
\]

But the graph is not zero.  The function does not vanish.  The integral is zero because the
negative and positive contributions balance exactly.

\begin{quote}
\textbf{An integral can be zero for two different reasons.}

It may be zero because there is no accumulation.

It may also be zero because positive and negative accumulation cancel.
\end{quote}

This distinction is crucial in motion.

If a car travels \(40\) miles east and then \(40\) miles west, its displacement is \(0\).  But
it did not stay still.  It traveled \(80\) miles.

The signed integral measures net change.  It does not measure total activity.

\subsection{Velocity: displacement versus distance traveled}
\label{subsec:velocity-displacement-distance-integrals}

Suppose the velocity of a particle is

\[
v(t)=
\begin{cases}
2, & 0\le t<1,\\
-1, & 1\le t<3,\\
3, & 3\le t\le4.
\end{cases}
\]

The signed displacement is

\[
\int_0^4 v(t)\,dt.
\]

Since the velocity is constant on each piece, we can compute this with rectangles:

\[
\int_0^4 v(t)\,dt
=
2(1)+(-1)(2)+3(1)
=
2-2+3
=
3.
\]

So the particle ends \(3\) units ahead of where it started.

\begin{figure}[h]
\centering
\begin{tikzpicture}
\begin{axis}[
    width=0.78\textwidth,
    height=0.48\textwidth,
    xmin=0, xmax=4.35,
    ymin=-1.6, ymax=3.5,
    axis lines=middle,
    xlabel={\(t\)},
    ylabel={\(v(t)\)},
    xtick={0,1,3,4},
    ytick={-1,0,2,3},
    grid=both,
    grid style={line width=.1pt, draw=gray!25},
]
\foreach \a/\b/\height in {0/1/2,1/3/-1,3/4/3} {
    \edef\temp{
        \noexpand\draw[fill=gray!15, draw=gray!60, fill opacity=0.5]
            (axis cs:\a,0) rectangle (axis cs:\b,\height);
    }
    \temp
}
\addplot[very thick] coordinates {(0,2) (1,2)};
\addplot[very thick] coordinates {(1,-1) (3,-1)};
\addplot[very thick] coordinates {(3,3) (4,3)};
\draw[dashed] (axis cs:1,-1.45) -- (axis cs:1,3.3);
\draw[dashed] (axis cs:3,-1.45) -- (axis cs:3,3.3);
\node[anchor=south] at (axis cs:0.5,2) {\(+2\)};
\node[anchor=north] at (axis cs:2,-1) {\(-2\)};
\node[anchor=south] at (axis cs:3.5,3) {\(+3\)};
\end{axis}
\end{tikzpicture}
\caption{The signed rectangle areas are \(2\), \(-2\), and \(3\).  Their sum is the displacement.}
\label{fig:piecewise-velocity-signed-area}
\end{figure}

The total distance traveled is different.  It counts backward motion positively:

\[
2(1)+1(2)+3(1)
=
2+2+3
=
7.
\]

The formula for total distance is

\[
\boxed{
\text{distance traveled}
=
\int_a^b |v(t)|\,dt.
}
\]

The absolute value removes the sign of the velocity before accumulation.

\[
\boxed{
\text{displacement}
=
\int_a^b v(t)\,dt.
}
\]

The signed integral answers:

\[
\text{Where did the object end compared with where it started?}
\]

The integral of the absolute value answers:

\[
\text{How much ground did the object cover?}
\]

\subsection{Total accumulation from absolute value}
\label{subsec:total-accumulation-absolute-value}

The same distinction appears outside motion.

Suppose \(f(x)\) represents a signed rate of change.  Positive values increase a quantity.
Negative values decrease it.

Then

\[
\int_a^b f(x)\,dx
\]

measures the \emph{net change}.  Positive and negative contributions may cancel.

But

\[
\int_a^b |f(x)|\,dx
\]

measures the \emph{total size of the accumulated change}, ignoring direction.

For the earlier example \(f(x)=x-1\) on \([0,3]\), we found

\[
\int_0^3 (x-1)\,dx=\frac32.
\]

The total accumulated size is

\[
\int_0^3 |x-1|\,dx.
\]

This adds both triangles positively:

\[
\int_0^3 |x-1|\,dx
=
\frac12+2
=
\frac52.
\]

So

\[
\text{net signed accumulation}=\frac32,
\]

while

\[
\text{total accumulation}=\frac52.
\]

The total is larger because cancellation has been removed.

\begin{quote}
\textbf{Net versus total.}

\[
\boxed{
\int_a^b f(x)\,dx
=
\text{positive accumulation}
-
\text{negative accumulation}.
}
\]

\[
\boxed{
\int_a^b |f(x)|\,dx
=
\text{positive accumulation}
+
\text{negative accumulation}.
}
\]
\end{quote}

The definite integral is therefore more flexible than ordinary area.  It can measure area,
displacement, accumulated signed change, or cancellation.  To use it well, we must always
ask what the sign means.

The next section collects the algebraic rules that make integrals useful.  These rules will
look familiar because they come from the same place as the integral itself: finite sums.

\section{Properties of the integral}
\label{sec:properties-of-the-integral}

The integral was born from sums.  Its basic properties are inherited from sums.

A Riemann sum has the form

\[
\sum_{j=1}^{n} f(x_j^*)\Delta x_j.
\]

Finite sums are linear.  They can be split into pieces.  They preserve inequalities.  They
respect symmetry.  When the Riemann sums approach a limit, those same habits survive
in the integral.

This section is not a list of tricks.  It is the arithmetic of accumulation.

\subsection{Linearity}
\label{subsec:integral-linearity}

Suppose two functions \(f\) and \(g\) are integrable on \([a,b]\).  If we accumulate their
sum, we get the sum of their accumulations:

\[
\int_a^b \bigl(f(x)+g(x)\bigr)\,dx
=
\int_a^b f(x)\,dx+\int_a^b g(x)\,dx.
\]

If we multiply a function by a constant \(c\), the accumulated amount is multiplied by
the same constant:

\[
\int_a^b c f(x)\,dx
=
c\int_a^b f(x)\,dx.
\]

Together these give the linearity rule.

\begin{quote}
\textbf{Linearity of the integral.}
If \(f\) and \(g\) are integrable on \([a,b]\), and \(A\) and \(B\) are constants, then

\[
\boxed{
\int_a^b \bigl(Af(x)+Bg(x)\bigr)\,dx
=
A\int_a^b f(x)\,dx
+
B\int_a^b g(x)\,dx.
}
\]
\end{quote}

\paragraph{Why it is true.}
For one partition and one set of sample points,

\[
\sum_{j=1}^{n} \bigl(Af(x_j^*)+Bg(x_j^*)\bigr)\Delta x_j
\]

equals

\[
A\sum_{j=1}^{n} f(x_j^*)\Delta x_j
+
B\sum_{j=1}^{n} g(x_j^*)\Delta x_j.
\]

This is just the distributive law for finite sums.  As the widths of the subintervals go to
zero, the left side approaches

\[
\int_a^b (Af+Bg)\,dx,
\]

and the two sums on the right approach

\[
A\int_a^b f\,dx
+
B\int_a^b g\,dx.
\]

So the integral obeys the same linearity.

\paragraph{Example.}
Find

\[
\int_0^4 (3+2x)\,dx
\]

using geometry.

By linearity,

\[
\int_0^4 (3+2x)\,dx
=
\int_0^4 3\,dx+2\int_0^4 x\,dx.
\]

The first integral is a rectangle:

\[
\int_0^4 3\,dx=12.
\]

The second is the area of a triangle with base \(4\) and height \(4\):

\[
\int_0^4 x\,dx=\frac12(4)(4)=8.
\]

Therefore

\[
\int_0^4 (3+2x)\,dx
=
12+2(8)
=
28.
\]

Linearity lets us break a complicated accumulation into simpler accumulations.

\subsection{Additivity over intervals}
\label{subsec:integral-additivity-intervals}

Accumulation over a long interval can be split into accumulation over shorter intervals.

If \(a<c<b\), then

\[
\int_a^b f(x)\,dx
\]

accumulates from \(a\) to \(b\).  We may first accumulate from \(a\) to \(c\), then from
\(c\) to \(b\):

\[
\boxed{
\int_a^b f(x)\,dx
=
\int_a^c f(x)\,dx
+
\int_c^b f(x)\,dx.
}
\]

\begin{figure}[h]
\centering
\begin{tikzpicture}
\begin{axis}[
    width=0.78\textwidth,
    height=0.48\textwidth,
    xmin=0, xmax=5.2,
    ymin=0, ymax=5.2,
    axis lines=left,
    xlabel={\(x\)},
    ylabel={\(y\)},
    xtick={0,2,5},
    xticklabels={\(a\),\(c\),\(b\)},
    ytick={0,1,2,3,4,5},
    grid=both,
    grid style={line width=.1pt, draw=gray!25},
]
\addplot[domain=0:2, samples=100, fill=gray!15, draw=none] {1+0.25*x+0.10*x^2} \closedcycle;
\addplot[domain=2:5, samples=100, fill=gray!25, draw=none] {1+0.25*x+0.10*x^2} \closedcycle;
\addplot[very thick, domain=0:5, samples=150] {1+0.25*x+0.10*x^2};
\draw[dashed] (axis cs:2,0) -- (axis cs:2,1.9);
\node[anchor=west] at (axis cs:0.45,0.55) {\(\int_a^c f\)};
\node[anchor=west] at (axis cs:3.0,0.85) {\(\int_c^b f\)};
\node[anchor=west] at (axis cs:2.7,3.2) {\(y=f(x)\)};
\end{axis}
\end{tikzpicture}
\caption{Accumulation over \([a,b]\) is the sum of accumulation over \([a,c]\) and \([c,b]\).}
\label{fig:additivity-over-intervals}
\end{figure}

\paragraph{Example.}
Suppose a particle has velocity \(v(t)\).  Then

\[
\int_0^5 v(t)\,dt
=
\int_0^2 v(t)\,dt+\int_2^5 v(t)\,dt.
\]

The displacement over \(5\) seconds is the displacement over the first \(2\) seconds plus
the displacement over the next \(3\) seconds.  This remains true even if one of the pieces
is negative.

\subsection{Orientation of the interval}
\label{subsec:orientation-of-interval-integral}

So far, we have usually assumed \(a<b\).  But the notation

\[
\int_a^b f(x)\,dx
\]

also remembers direction.

Moving from \(a\) to \(b\) is the opposite of moving from \(b\) to \(a\).  To make the
algebra of accumulation respect that direction, we define

\[
\boxed{
\int_b^a f(x)\,dx=-\int_a^b f(x)\,dx.
}
\]

Also,

\[
\boxed{
\int_a^a f(x)\,dx=0.
}
\]

There is no width, so there is no accumulation.

With these conventions, the interval-addition rule becomes especially clean:

\[
\boxed{
\int_a^b f(x)\,dx+\int_b^c f(x)\,dx=\int_a^c f(x)\,dx.
}
\]

This formula works no matter where \(a\), \(b\), and \(c\) lie on the number line, as long
as the integrals exist.

\paragraph{Warning.}
Order matters when the interval is reversed.  Even though

\[
1>0,
\]

we have

\[
\int_1^0 1\,dx=-1.
\]

The function is positive, but the direction of travel is negative.  The symbol \(dx\) carries
that orientation.

For order properties, unless we say otherwise, we assume

\[
a<b.
\]

\subsection{Order properties}
\label{subsec:integral-order-properties}

If one function is always below another, its integral is no larger.

\begin{quote}
\textbf{Order property.}
Assume \(a<b\).  If \(f\) and \(g\) are integrable on \([a,b]\), and

\[
f(x)\le g(x)
\]

for every \(x\in[a,b]\), then

\[
\boxed{
\int_a^b f(x)\,dx
\le
\int_a^b g(x)\,dx.
}
\]
\end{quote}

\paragraph{Why it is true.}
For every sample point,

\[
f(x_j^*)\le g(x_j^*).
\]

Since \(\Delta x_j>0\),

\[
f(x_j^*)\Delta x_j\le g(x_j^*)\Delta x_j.
\]

Adding gives

\[
\sum_{j=1}^{n} f(x_j^*)\Delta x_j
\le
\sum_{j=1}^{n} g(x_j^*)\Delta x_j.
\]

Taking limits preserves the inequality.

Several useful facts follow immediately.

If

\[
f(x)\ge0
\]

on \([a,b]\), then

\[
\boxed{
\int_a^b f(x)\,dx\ge0.
}
\]

If

\[
m\le f(x)\le M
\]

on \([a,b]\), then

\[
\boxed{
m(b-a)
\le
\int_a^b f(x)\,dx
\le
M(b-a).
}
\]

This says that the accumulation of \(f\) is trapped between two rectangles: one with
height \(m\), and one with height \(M\).

\paragraph{Example.}
Suppose a machine uses power \(P(t)\) between \(2\) and \(5\) kilowatts during a
\(6\)-hour work period.  The total energy used is

\[
\int_0^6 P(t)\,dt.
\]

Because

\[
2\le P(t)\le5,
\]

we know

\[
2(6)\le \int_0^6 P(t)\,dt\le5(6).
\]

So the energy used is between

\[
12
\]

and

\[
30
\]

kilowatt-hours.

We do not need the exact power curve to get a useful bound.

\subsection{Net change is no larger than total change}
\label{subsec:net-change-total-change-inequality}

The order property also explains the relation between signed accumulation and total
accumulation.

For every \(x\),

\[
-|f(x)|\le f(x)\le |f(x)|.
\]

Integrating over \([a,b]\), with \(a<b\), gives

\[
-\int_a^b |f(x)|\,dx
\le
\int_a^b f(x)\,dx
\le
\int_a^b |f(x)|\,dx.
\]

This is the same as

\[
\boxed{
\left|\int_a^b f(x)\,dx\right|
\le
\int_a^b |f(x)|\,dx.
}
\]

In words:

\[
\boxed{
\text{the size of the net accumulation is at most the total accumulation.}
}
\]

This matches common sense.  A person who walks forward and backward may end close to
the starting point, but the total walking distance cannot be smaller than the net
displacement.

Equality happens only when there is no cancellation, or when all cancellation is absent in
one direction.  For instance, if \(f(x)\ge0\) on the whole interval, then

\[
\left|\int_a^b f(x)\,dx\right|
=
\int_a^b |f(x)|\,dx.
\]

If \(f\) changes sign, the inequality is usually strict.

\subsection{Average value of a function}
\label{subsec:average-value-function}

A definite integral accumulates total signed amount.  Sometimes we want to spread that
amount evenly over the interval.

Suppose \(f(x)\) represents temperature over time.  The integral

\[
\int_a^b f(x)\,dx
\]

has units

\[
\text{temperature}\cdot\text{time}.
\]

To get an average temperature, divide by the time length \(b-a\).

\begin{quote}
\textbf{Average value.}
If \(f\) is integrable on \([a,b]\), with \(a<b\), then the average value of \(f\) on
\([a,b]\) is

\[
\boxed{
f_{\text{avg}}
=
\frac{1}{b-a}\int_a^b f(x)\,dx.
}
\]
\end{quote}

This number is the constant height that gives the same signed area:

\[
f_{\text{avg}}(b-a)=\int_a^b f(x)\,dx.
\]

\paragraph{Example.}
Find the average value of

\[
f(x)=1+x
\]

on \([0,4]\).

The region under the graph is a trapezoid.  Its parallel sides have lengths

\[
f(0)=1
\]

and

\[
f(4)=5,
\]

and its width is \(4\).  Therefore

\[
\int_0^4 (1+x)\,dx
=
\frac12(1+5)(4)
=
12.
\]

The average value is

\[
f_{\text{avg}}
=
\frac{1}{4}\cdot12
=
3.
\]

So the slanted graph \(y=1+x\) has the same accumulated area over \([0,4]\) as the
constant function \(y=3\).

\begin{figure}[h]
\centering
\begin{minipage}{0.47\textwidth}
\centering
\begin{tikzpicture}
\begin{axis}[
    width=\textwidth,
    height=0.72\textwidth,
    xmin=0, xmax=4.4,
    ymin=0, ymax=5.6,
    axis lines=left,
    xlabel={\(x\)},
    ylabel={\(y\)},
    xtick={0,4},
    ytick={1,3,5},
    grid=both,
    grid style={line width=.1pt, draw=gray!25},
]
\draw[fill=gray!15, draw=gray!60]
  (axis cs:0,0) -- (axis cs:0,1) -- (axis cs:4,5) -- (axis cs:4,0) -- cycle;
\addplot[very thick, domain=0:4, samples=2] {1+x};
\node[anchor=west] at (axis cs:1.1,4.2) {\(y=1+x\)};
\node at (axis cs:2,1.25) {area \(=12\)};
\end{axis}
\end{tikzpicture}

actual graph
\end{minipage}
\hfill
\begin{minipage}{0.47\textwidth}
\centering
\begin{tikzpicture}
\begin{axis}[
    width=\textwidth,
    height=0.72\textwidth,
    xmin=0, xmax=4.4,
    ymin=0, ymax=5.6,
    axis lines=left,
    xlabel={\(x\)},
    ylabel={\(y\)},
    xtick={0,4},
    ytick={1,3,5},
    grid=both,
    grid style={line width=.1pt, draw=gray!25},
]
\draw[fill=gray!15, draw=gray!60]
  (axis cs:0,0) rectangle (axis cs:4,3);
\addplot[very thick, domain=0:4, samples=2] {3};
\node[anchor=west] at (axis cs:1.1,3.25) {\(y=3\)};
\node at (axis cs:2,1.25) {area \(=12\)};
\end{axis}
\end{tikzpicture}

average height
\end{minipage}
\caption{The average value is the constant height that gives the same accumulated area.}
\label{fig:average-value-equal-area}
\end{figure}

\subsection{The average value is really attained when \(f\) is continuous}
\label{subsec:average-value-attained}

For a continuous function, the average value is not only a number.  It is a height the
function actually reaches.

\begin{quote}
\textbf{Mean Value Theorem for Integrals.}
If \(f\) is continuous on \([a,b]\), with \(a<b\), then there is some point
\(c\in[a,b]\) such that

\[
\boxed{
f(c)=\frac{1}{b-a}\int_a^b f(x)\,dx.
}
\]
\end{quote}

\paragraph{Why it is true.}
Because \(f\) is continuous on a closed interval, it has a minimum value \(m\) and a
maximum value \(M\).  Thus

\[
m\le f(x)\le M
\]

for every \(x\in[a,b]\).  The order property gives

\[
m(b-a)
\le
\int_a^b f(x)\,dx
\le
M(b-a).
\]

Divide by \(b-a>0\):

\[
m
\le
\frac{1}{b-a}\int_a^b f(x)\,dx
\le
M.
\]

So the average value lies between the minimum and maximum values of \(f\).  Since
\(f\) is continuous, the Intermediate Value Theorem says that \(f\) takes every value
between \(m\) and \(M\).  Therefore it takes the average value.

\paragraph{Why continuity matters.}
Consider the step function

\[
f(x)=
\begin{cases}
0, & 0\le x\le \frac12,\\
1, & \frac12<x\le1.
\end{cases}
\]

The area under the graph is

\[
\int_0^1 f(x)\,dx=\frac12.
\]

So the average value is

\[
f_{\text{avg}}=\frac12.
\]

But the function only takes the values \(0\) and \(1\).  There is no point \(c\) where

\[
f(c)=\frac12.
\]

The function is integrable, but it is not continuous.  The average value exists, but the
function never actually equals its average.

This is why the theorem requires continuity.

\subsection{Symmetry}
\label{subsec:integral-symmetry}

Symmetry can make integrals much easier.

Recall two kinds of symmetry:

\[
f(-x)=f(x)
\]

means \(f\) is even.  Its graph is symmetric across the \(y\)-axis.

\[
f(-x)=-f(x)
\]

means \(f\) is odd.  Its graph is symmetric through the origin.

On a symmetric interval \([-a,a]\), these symmetries become integral shortcuts.

\begin{quote}
\textbf{Symmetry rules.}
If \(f\) is integrable on \([-a,a]\), then:

\[
\boxed{
\text{if }f\text{ is even,}\qquad
\int_{-a}^{a} f(x)\,dx
=
2\int_0^a f(x)\,dx.
}
\]

\[
\boxed{
\text{if }f\text{ is odd,}\qquad
\int_{-a}^{a} f(x)\,dx
=
0.
}
\]
\end{quote}

For an even function, the left half and right half contribute equal signed areas.

For an odd function, the left half and right half contribute opposite signed areas.  They
cancel.

\begin{figure}[h]
\centering
\begin{tikzpicture}
\begin{axis}[
    width=0.78\textwidth,
    height=0.48\textwidth,
    xmin=-1.35, xmax=1.35,
    ymin=-1.35, ymax=1.35,
    axis lines=middle,
    axis on top,
    xlabel={\(x\)},
    ylabel={\(y\)},
    xtick={-1,0,1},
    ytick={-1,0,1},
    grid=both,
    grid style={line width=.1pt, draw=gray!25},
]
\addplot[domain=-1:0, samples=100, fill=gray!25, draw=none] {x^3} \closedcycle;
\addplot[domain=0:1, samples=100, fill=gray!15, draw=none, fill opacity=0.5] {x^3} \closedcycle;
\addplot[very thick, domain=-1.15:1.15, samples=150] {x^3};
\node[anchor=south east] at (axis cs:-0.25,-0.55) {negative};
\node[anchor=south west] at (axis cs:0.35,0.25) {positive};
\node[anchor=west] at (axis cs:0.55,0.9) {\(y=x^3\)};
\end{axis}
\end{tikzpicture}
\caption{For an odd function on a symmetric interval, the two signed areas cancel.}
\label{fig:odd-function-cancellation}
\end{figure}

\paragraph{Example.}
Since

\[
f(x)=x^3
\]

is odd,

\[
\int_{-2}^{2} x^3\,dx=0.
\]

No antiderivative is needed.  The cancellation comes from symmetry.

\paragraph{Example.}
The function

\[
g(x)=x^2+3
\]

is even.  Therefore

\[
\int_{-2}^{2} (x^2+3)\,dx
=
2\int_0^2 (x^2+3)\,dx.
\]

This does not yet give the final numerical value, but it cuts the work in half.  Later, the
Fundamental Theorem will make the remaining integral easy to evaluate.

\paragraph{Even plus odd.}
Many functions split into an even part and an odd part.  For example,

\[
x^2+x
\]

is the sum of the even function \(x^2\) and the odd function \(x\).  On \([-1,1]\),

\[
\int_{-1}^{1} (x^2+x)\,dx
=
\int_{-1}^{1} x^2\,dx+\int_{-1}^{1} x\,dx.
\]

The odd part cancels:

\[
\int_{-1}^{1} x\,dx=0.
\]

So

\[
\int_{-1}^{1} (x^2+x)\,dx
=
\int_{-1}^{1} x^2\,dx.
\]

Symmetry lets us see cancellation before doing any calculation.

\subsection{Where the properties come from}
\label{subsec:where-integral-properties-come-from}

The properties of the integral are not mysterious.  They are the properties of finite sums
surviving a limiting process.

\[
\begin{array}{c|c}
\textbf{finite sums} & \textbf{integrals} \\ \hline
\sum (Af_j+Bg_j)\Delta x_j
=
A\sum f_j\Delta x_j+B\sum g_j\Delta x_j
&
\int_a^b (Af+Bg)\,dx
=
A\int_a^b f\,dx+B\int_a^b g\,dx
\\[2mm]
\text{a long sum can be split into shorter sums}
&
\int_a^b f\,dx
=
\int_a^c f\,dx+\int_c^b f\,dx
\\[2mm]
f_j\le g_j \text{ for all }j
\Rightarrow
\sum f_j\Delta x_j\le \sum g_j\Delta x_j
&
f\le g
\Rightarrow
\int_a^b f\,dx\le \int_a^b g\,dx
\\[2mm]
\text{matched positive and negative terms can cancel}
&
\text{odd functions cancel on symmetric intervals}
\end{array}
\]

These rules make the integral usable before we know many exact antiderivatives.  They
allow us to split, estimate, compare, and simplify accumulated quantities.

One question remains before we can rely on the integral fully:

\[
\text{Which functions actually have integrals?}
\]

Continuous functions do.  So do many functions with jumps.  But not every bounded
function is integrable.  The next section explains where the line is.
\section{Existence of integrals}
\label{sec:existence-of-integrals}

The last section used the integral as if it were already available.  We split it, compared
it, averaged it, and used symmetry to simplify it.

But one question is still waiting:

\[
\text{When does a function actually have a Riemann integral?}
\]

For the functions that appear in most applications, the answer is reassuring.  Continuous
functions are integrable.  Monotone functions are integrable.  Functions with finitely
many jumps are integrable.

But boundedness alone is not enough.  A function can be bounded and still be too
wild for its Riemann sums to settle down.

The test is this:

\[
\text{Can upper and lower estimates be forced together?}
\]

\subsection{The upper-lower test}
\label{subsec:upper-lower-test-integrability}

Let \(f\) be bounded on \([a,b]\).  Choose a partition

\[
P:\qquad a=x_0<x_1<\cdots<x_n=b.
\]

On the \(j\)-th subinterval

\[
[x_{j-1},x_j],
\]

let

\[
m_j=\inf\{f(x):x_{j-1}\le x\le x_j\}
\]

and

\[
M_j=\sup\{f(x):x_{j-1}\le x\le x_j\}.
\]

For friendly functions, these are just the minimum and maximum values of \(f\) on
that subinterval.

The lower sum is

\[
L(P,f)=\sum_{j=1}^{n}m_j\Delta x_j,
\qquad
\Delta x_j=x_j-x_{j-1}.
\]

The upper sum is

\[
U(P,f)=\sum_{j=1}^{n}M_j\Delta x_j.
\]

Every Riemann sum using sample points inside these subintervals is trapped between
them:

\[
L(P,f)\le
\sum_{j=1}^{n}f(x_j^*)\Delta x_j
\le U(P,f).
\]

The gap is

\[
U(P,f)-L(P,f)
=
\sum_{j=1}^{n}(M_j-m_j)\Delta x_j.
\]

This gap measures the total uncertainty left by the partition.  If the graph does not
vary much inside each subinterval, then \(M_j-m_j\) is small.  If the graph jumps
or oscillates badly inside every subinterval, then the gap may refuse to shrink.

\begin{quote}
\textbf{Upper-lower integrability test.}

A bounded function \(f\) on \([a,b]\) is Riemann integrable if and only if, for every
\(\varepsilon>0\), there is a partition \(P\) such that

\[
\boxed{
U(P,f)-L(P,f)<\varepsilon.
}
\]

In words: the upper and lower estimates can be made as close as we wish.
\end{quote}

This criterion is the cleanest way to understand which functions are integrable.  It says
that the integral exists when the remaining uncertainty can be squeezed down to zero.

\subsection{Continuous functions are integrable}
\label{subsec:continuous-functions-integrable}

A continuous graph on a closed interval cannot jump.  If we look at a short enough
piece of the graph, the output values stay close together.

That is exactly what the upper-lower test needs.

\paragraph{A first example.}
Consider

\[
f(x)=x^2
\]

on \([0,1]\).  Divide the interval into \(n\) equal pieces:

\[
0,\frac1n,\frac2n,\ldots,\frac{n-1}{n},1.
\]

Since \(x^2\) is increasing on \([0,1]\), the lower height on

\[
\left[\frac{j-1}{n},\frac{j}{n}\right]
\]

is

\[
\left(\frac{j-1}{n}\right)^2,
\]

and the upper height is

\[
\left(\frac{j}{n}\right)^2.
\]

The width is \(1/n\).  Therefore

\[
U_n-L_n
=
\sum_{j=1}^{n}
\left[
\left(\frac{j}{n}\right)^2
-
\left(\frac{j-1}{n}\right)^2
\right]\frac1n.
\]

Factor out \(1/n\):

\[
U_n-L_n
=
\frac1n
\sum_{j=1}^{n}
\left[
\left(\frac{j}{n}\right)^2
-
\left(\frac{j-1}{n}\right)^2
\right].
\]

The sum telescopes:

\[
\sum_{j=1}^{n}
\left[
\left(\frac{j}{n}\right)^2
-
\left(\frac{j-1}{n}\right)^2
\right]
=
1^2-0^2
=
1.
\]

So

\[
U_n-L_n=\frac1n.
\]

As \(n\to\infty\),

\[
\frac1n\to0.
\]

Thus the upper and lower sums squeeze together, and \(x^2\) is integrable on
\([0,1]\).

The same idea works for every continuous function, although the arithmetic may not
telescope so neatly.

\begin{quote}
\textbf{Theorem.}
If \(f\) is continuous on a closed interval \([a,b]\), then \(f\) is Riemann integrable
on \([a,b]\).
\end{quote}

\paragraph{Proof idea.}
A continuous function on a closed interval has a special kind of steadiness:

\[
\text{if }x\text{ and }y\text{ are close enough, then }f(x)\text{ and }f(y)\text{ are close,}
\]

and the same meaning of ``close enough'' works everywhere on the interval.

More precisely, for every \(\eta>0\), there is a \(\delta>0\) such that

\[
|x-y|<\delta
\qquad\Longrightarrow\qquad
|f(x)-f(y)|<\eta
\]

for all \(x,y\in[a,b]\).

Now choose a partition whose subinterval widths are all less than \(\delta\).  Then on
each subinterval,

\[
M_j-m_j<\eta.
\]

Therefore

\[
U(P,f)-L(P,f)
=
\sum_{j=1}^{n}(M_j-m_j)\Delta x_j
<
\sum_{j=1}^{n}\eta\,\Delta x_j.
\]

But

\[
\sum_{j=1}^{n}\Delta x_j=b-a.
\]

So

\[
U(P,f)-L(P,f)<\eta(b-a).
\]

Given any desired error \(\varepsilon>0\), choose

\[
\eta=\frac{\varepsilon}{b-a}.
\]

Then

\[
U(P,f)-L(P,f)<\varepsilon.
\]

By the upper-lower test, \(f\) is integrable.

\paragraph{Why the closed interval matters.}
The function

\[
f(x)=\frac1x
\]

is continuous on \((0,1]\), but it is not bounded near \(0\).  It is not an ordinary
Riemann integrable function on \([0,1]\).

This does not mean we can never discuss

\[
\int_0^1 \frac1x\,dx.
\]

It means that such a question is not an ordinary definite integral.  It belongs to the
later topic of \emph{improper integrals}, where unbounded behavior must be handled
with a new limiting process.

\subsection{Monotone functions are integrable}
\label{subsec:monotone-functions-integrable}

Continuity is enough for integrability, but it is not necessary.

A function may jump and still be integrable.  The simplest reason is monotonicity:
a monotone function can move only one way.  It may jump, but it cannot oscillate
up and down forever inside every small interval.

Suppose \(f\) is increasing on \([a,b]\).  Divide \([a,b]\) into \(n\) equal pieces, each
of width

\[
\Delta x=\frac{b-a}{n}.
\]

On the \(j\)-th interval,

\[
[x_{j-1},x_j],
\]

the lower height is \(f(x_{j-1})\), and the upper height is \(f(x_j)\).  Therefore

\[
U_n-L_n
=
\sum_{j=1}^{n}\bigl(f(x_j)-f(x_{j-1})\bigr)\Delta x.
\]

Since \(\Delta x\) is the same for every interval,

\[
U_n-L_n
=
\Delta x
\sum_{j=1}^{n}\bigl(f(x_j)-f(x_{j-1})\bigr).
\]

The sum telescopes:

\[
\sum_{j=1}^{n}\bigl(f(x_j)-f(x_{j-1})\bigr)
=
f(b)-f(a).
\]

Thus

\[
U_n-L_n
=
\frac{b-a}{n}\bigl(f(b)-f(a)\bigr).
\]

As \(n\to\infty\), this goes to \(0\).  So the upper and lower sums squeeze together.

\begin{quote}
\textbf{Theorem.}
If \(f\) is monotone on \([a,b]\), then \(f\) is Riemann integrable on \([a,b]\).
\end{quote}

For a decreasing function, the same proof works with \(f(a)-f(b)\) in place of
\(f(b)-f(a)\).

\begin{figure}[h]
\centering
\begin{minipage}{0.47\textwidth}
\centering
\begin{tikzpicture}
\begin{axis}[
    width=\textwidth,
    height=0.72\textwidth,
    xmin=0, xmax=4.25,
    ymin=0, ymax=3.35,
    axis lines=left,
    xlabel={\(x\)},
    ylabel={\(y\)},
    xtick={0,1,2,3,4},
    ytick={0,1,2,3},
    grid=both,
    grid style={line width=.1pt, draw=gray!25},
]
\foreach \xleft/\xright in {0/1,1/2,2/3,3/4} {
    \pgfmathsetmacro{\rectheight}{1+0.20*\xleft+0.12*\xleft*\xleft}
    \edef\temp{
        \noexpand\draw[fill=gray!15, draw=gray!60]
            (axis cs:\xleft,0) rectangle (axis cs:\xright,\rectheight);
    }
    \temp
}
\addplot[very thick, domain=0:4, samples=150] {1+0.20*x+0.12*x^2};
\node[anchor=west] at (axis cs:0.25,2.75) {lower sum};
\end{axis}
\end{tikzpicture}
\end{minipage}
\hfill
\begin{minipage}{0.47\textwidth}
\centering
\begin{tikzpicture}
\begin{axis}[
    width=\textwidth,
    height=0.72\textwidth,
    xmin=0, xmax=4.25,
    ymin=0, ymax=3.35,
    axis lines=left,
    xlabel={\(x\)},
    ylabel={\(y\)},
    xtick={0,1,2,3,4},
    ytick={0,1,2,3},
    grid=both,
    grid style={line width=.1pt, draw=gray!25},
]
\foreach \xleft/\xright in {0/1,1/2,2/3,3/4} {
    \pgfmathsetmacro{\rectheight}{1+0.20*\xright+0.12*\xright*\xright}
    \edef\temp{
        \noexpand\draw[fill=gray!15, draw=gray!60]
            (axis cs:\xleft,0) rectangle (axis cs:\xright,\rectheight);
    }
    \temp
}
\addplot[very thick, domain=0:4, samples=150] {1+0.20*x+0.12*x^2};
\node[anchor=west] at (axis cs:0.25,2.75) {upper sum};
\end{axis}
\end{tikzpicture}
\end{minipage}
\caption{For an increasing function, left endpoint rectangles give lower sums and right endpoint rectangles give upper sums.  As the widths shrink, the two sums squeeze together.}
\label{fig:monotone-upper-lower-sums}
\end{figure}

This theorem explains why jumps are not automatically dangerous.  A step function
can be monotone, and a monotone step function is integrable.

\paragraph{Example.}
Let

\[
H(x)=
\begin{cases}
0, & 0\le x<\frac12,\\
1, & \frac12\le x\le1.
\end{cases}
\]

This function jumps at \(x=\frac12\), but it is increasing.  Therefore it is integrable.
The area under the graph is a rectangle of width \(1/2\) and height \(1\), so

\[
\int_0^1 H(x)\,dx=\frac12.
\]

The value at the single point \(x=\frac12\) does not affect the integral.  A single point
has no width.

\subsection{Functions with finitely many jump discontinuities}
\label{subsec:finitely-many-jumps-integrable}

A function may have several jumps and still be integrable.  The reason is the same:
finitely many points can be covered by intervals whose total width is as small as we
wish.

\begin{quote}
\textbf{Theorem.}
If \(f\) is bounded on \([a,b]\) and is continuous except at finitely many points, then
\(f\) is Riemann integrable on \([a,b]\).
\end{quote}

\paragraph{Proof idea.}
Suppose the discontinuities occur at finitely many points.  Put tiny intervals around
those points.  Since there are only finitely many of them, we can make the total width
of all the tiny intervals very small.

On those tiny intervals, the function may behave badly, but it is bounded.  So the
upper-lower contribution is at worst

\[
\text{bounded height}\times\text{tiny total width},
\]

which can be made small.

On the rest of the interval, the function is continuous on closed pieces.  So the
upper-lower contribution there can also be made small.

Putting the two parts together, the total upper-lower gap can be made smaller than
any chosen \(\varepsilon\).  Hence the function is integrable.

\paragraph{Example.}
Consider

\[
f(x)=
\begin{cases}
x, & 0\le x<1,\\
x+2, & 1\le x\le2.
\end{cases}
\]

The graph has one jump at \(x=1\).

\begin{figure}[h]
\centering
\begin{tikzpicture}
\begin{axis}[
    width=0.78\textwidth,
    height=0.48\textwidth,
    xmin=0, xmax=2.25,
    ymin=0, ymax=4.4,
    axis lines=left,
    xlabel={\(x\)},
    ylabel={\(y\)},
    xtick={0,1,2},
    ytick={0,1,3,4},
    grid=both,
    grid style={line width=.1pt, draw=gray!25},
]
\draw[fill=gray!15, draw=none]
    (axis cs:0,0) -- (axis cs:1,1) -- (axis cs:1,0) -- cycle;
\draw[fill=gray!25, draw=none]
    (axis cs:1,0) -- (axis cs:1,3) -- (axis cs:2,4) -- (axis cs:2,0) -- cycle;
\addplot[very thick, domain=0:1, samples=2] {x};
\addplot[very thick, domain=1:2, samples=2] {x+2};
\addplot[only marks, mark=o] coordinates {(1,1)};
\addplot[only marks, mark=*] coordinates {(1,3)};
\draw[dashed] (axis cs:1,0) -- (axis cs:1,3.2);
\node[anchor=west] at (axis cs:0.25,0.65) {\(y=x\)};
\node[anchor=west] at (axis cs:1.25,3.85) {\(y=x+2\)};
\node[anchor=west] at (axis cs:1.04,1.9) {jump};
\end{axis}
\end{tikzpicture}
\caption{A finite jump does not prevent integrability.  The jump occurs at one point, and one point has no width.}
\label{fig:finite-jump-integrable}
\end{figure}

The integral exists.  We can compute it by ordinary geometry.

On \([0,1]\), the area under \(y=x\) is a triangle:

\[
\frac12.
\]

On \([1,2]\), the area under \(y=x+2\) is a trapezoid of width \(1\), with heights
\(3\) and \(4\).  Its area is

\[
\frac12(3+4)(1)=\frac72.
\]

Therefore

\[
\int_0^2 f(x)\,dx
=
\frac12+\frac72
=
4.
\]

If we changed the value of \(f(1)\), the integral would still be \(4\).  The height at one
point cannot change accumulated area.

\begin{quote}
\textbf{Changing finitely many values does not change a Riemann integral.}

If two bounded functions differ only at finitely many points, then either both are
integrable with the same integral, or neither is integrable.
\end{quote}

\subsection{A bounded function that is not Riemann integrable}
\label{subsec:bounded-not-riemann-integrable}

Now for the warning.

Define

\[
D(x)=
\begin{cases}
1, & x\text{ is rational},\\
0, & x\text{ is irrational}
\end{cases}
\]

on \([0,1]\).

This function is bounded:

\[
0\le D(x)\le1.
\]

But it is not Riemann integrable.

The reason is that every interval, no matter how short, contains rational numbers and
irrational numbers.  Therefore, on every subinterval,

\[
m_j=0
\]

and

\[
M_j=1.
\]

So for every partition \(P\),

\[
L(P,D)=0
\]

and

\[
U(P,D)=1.
\]

The upper-lower gap is always

\[
U(P,D)-L(P,D)=1.
\]

It never shrinks.

\begin{figure}[h]
\centering
\begin{tikzpicture}
\begin{axis}[
    width=0.78\textwidth,
    height=0.42\textwidth,
    xmin=0, xmax=1,
    ymin=-0.25, ymax=1.25,
    axis lines=left,
    xlabel={\(x\)},
    ylabel={\(D(x)\)},
    xtick={0,0.25,0.5,0.75,1},
    ytick={0,1},
    grid=both,
    grid style={line width=.1pt, draw=gray!25},
]
\addplot[only marks, mark=*, mark size=1.4pt] coordinates {
    (0.03,1) (0.07,1) (0.11,1) (0.16,1) (0.22,1) (0.29,1)
    (0.34,1) (0.41,1) (0.48,1) (0.53,1) (0.61,1) (0.67,1)
    (0.72,1) (0.81,1) (0.88,1) (0.94,1)
};
\addplot[only marks, mark=*, mark size=1.4pt] coordinates {
    (0.02,0) (0.09,0) (0.14,0) (0.19,0) (0.27,0) (0.31,0)
    (0.38,0) (0.45,0) (0.57,0) (0.63,0) (0.70,0) (0.76,0)
    (0.84,0) (0.91,0) (0.97,0)
};
\node[anchor=west] at (axis cs:0.12,1.08) {rational inputs give height \(1\)};
\node[anchor=west] at (axis cs:0.12,0.13) {irrational inputs give height \(0\)};
\end{axis}
\end{tikzpicture}
\caption{This picture can only hint at the function \(D\).  In every interval, the graph has points at height \(1\) and points at height \(0\).}
\label{fig:dirichlet-function-not-integrable}
\end{figure}

This example is extreme.  It is discontinuous at every point.  It shows why the
definition of the integral cannot simply say ``bounded functions have area.''

They do not.

\begin{quote}
\textbf{Safe classes of integrable functions.}

The following functions are Riemann integrable on closed intervals:

\[
\begin{array}{c|c}
\textbf{kind of function} & \textbf{reason} \\ \hline
\text{continuous functions} & \text{small intervals have small oscillation} \\
\text{monotone functions} & \text{upper-lower gaps telescope} \\
\text{bounded functions with} & \text{discontinuity points can be covered} \\
\text{finitely many discontinuities} & \text{by tiny intervals}
\end{array}
\]
\paragraph{The general rule.}
The ultimate test for Riemann integrability was proven by Henri Lebesgue. A bounded function is integrable if and only if its set of discontinuity points has \emph{measure zero}---meaning all the bad points can be covered by a collection of intervals whose total cumulative length is arbitrarily small.
\end{quote}

The integral is now on solid ground.  We know what it means, how to estimate it,
what signs do, what algebraic rules it obeys, and which common functions actually
have one.

The chapter ends by practicing the setup.  The next chapter will reveal the surprise:
if the upper endpoint moves, an integral becomes a function, and its derivative is the
function we started with.

\section{Chapter review and discovery problems}
\label{sec:chapter8-review-discovery}

This chapter built the integral from the ground up.

We started with velocity and distance:

\[
\text{distance}\approx \sum \text{velocity}\cdot\text{time}.
\]

Then we saw the geometric version:

\[
\text{area}\approx \sum \text{height}\cdot\text{width}.
\]

The definite integral is the limiting version of those sums:

\[
\int_a^b f(x)\,dx.
\]

It measures signed accumulation.  When \(f\ge0\), it also measures ordinary area.
When \(f\) changes sign, it measures positive accumulation minus negative
accumulation.

\[
\begin{array}{c|c|c}
\textbf{idea} & \textbf{formula} & \textbf{meaning} \\ \hline
\text{Riemann sum}
&
\displaystyle \sum_{j=1}^n f(x_j^*)\Delta x_j
&
\text{finite estimate}
\\[2mm]
\text{definite integral}
&
\displaystyle \int_a^b f(x)\,dx
&
\text{limit of Riemann sums}
\\[2mm]
\text{displacement}
&
\displaystyle \int_a^b v(t)\,dt
&
\text{net change in position}
\\[2mm]
\text{distance traveled}
&
\displaystyle \int_a^b |v(t)|\,dt
&
\text{total motion}
\\[2mm]
\text{average value}
&
\displaystyle \frac{1}{b-a}\int_a^b f(x)\,dx
&
\text{constant height with same area}
\end{array}
\]

\subsection{Skills: set up Riemann sums}
\label{subsec:skills-set-up-riemann-sums}

A Riemann sum setup has four parts:

\[
\text{interval},\qquad
\text{partition},\qquad
\text{sample points},\qquad
\text{sum}.
\]

For an equal partition of \([a,b]\) into \(n\) pieces,

\[
\Delta x=\frac{b-a}{n}.
\]

The endpoints are

\[
x_j=a+j\Delta x.
\]

Right endpoints use \(x_j\).  Left endpoints use \(x_{j-1}\).  Midpoints use

\[
\frac{x_{j-1}+x_j}{2}
=
a+\left(j-\frac12\right)\Delta x.
\]

\paragraph{Worked example.}
Write

\[
\int_0^2 (1+x^2)\,dx
\]

as a limit of right endpoint Riemann sums.  Then evaluate the limit using the power
sum formula.

The interval is \([0,2]\), so

\[
\Delta x=\frac{2}{n}.
\]

The right endpoint of the \(j\)-th subinterval is

\[
x_j=\frac{2j}{n}.
\]

Thus

\[
\int_0^2 (1+x^2)\,dx
=
\lim_{n\to\infty}
\sum_{j=1}^{n}
\left(1+\left(\frac{2j}{n}\right)^2\right)\frac{2}{n}.
\]

Simplify the sum:

\[
\sum_{j=1}^{n}
\left(1+\left(\frac{2j}{n}\right)^2\right)\frac{2}{n}
=
\sum_{j=1}^{n}\frac{2}{n}
+
\sum_{j=1}^{n}\frac{8j^2}{n^3}.
\]

The first sum is

\[
\sum_{j=1}^{n}\frac{2}{n}=2.
\]

The second sum is

\[
\frac{8}{n^3}\sum_{j=1}^{n}j^2.
\]

Using

\[
\sum_{j=1}^{n}j^2=\frac{n(n+1)(2n+1)}{6},
\]

we get

\[
\frac{8}{n^3}\sum_{j=1}^{n}j^2
=
\frac{8}{n^3}\cdot \frac{n(n+1)(2n+1)}{6}.
\]

As \(n\to\infty\),

\[
\frac{n(n+1)(2n+1)}{n^3}\to2.
\]

Therefore

\[
\frac{8}{n^3}\cdot \frac{n(n+1)(2n+1)}{6}
\to
\frac{8\cdot2}{6}
=
\frac83.
\]

So

\[
\boxed{
\int_0^2 (1+x^2)\,dx
=
2+\frac83
=
\frac{14}{3}.
}
\]

No antiderivative was used.  This came directly from the definition.

\paragraph{Setup practice.}
For each integral, write a Riemann sum using the requested sample points.  Do not
evaluate the limit unless asked.

\begin{enumerate}
\item Use left endpoints:

\[
\int_1^5 \sqrt{x}\,dx.
\]

Here

\[
\Delta x=\frac4n,
\qquad
x_{j-1}=1+\frac{4(j-1)}{n}.
\]

So the sum is

\[
\sum_{j=1}^{n}
\sqrt{1+\frac{4(j-1)}{n}}\cdot \frac4n.
\]

\item Use midpoints:

\[
\int_{-2}^{2}(3-x^2)\,dx.
\]

Here

\[
\Delta x=\frac4n,
\qquad
x_j^*=-2+\left(j-\frac12\right)\frac4n.
\]

So the sum is

\[
\sum_{j=1}^{n}
\left[
3-
\left(
-2+\left(j-\frac12\right)\frac4n
\right)^2
\right]\frac4n.
\]

\item Use right endpoints:

\[
\int_0^\pi \sin x\,dx.
\]

Here

\[
\Delta x=\frac{\pi}{n},
\qquad
x_j=\frac{j\pi}{n}.
\]

So the sum is

\[
\sum_{j=1}^{n}
\sin\left(\frac{j\pi}{n}\right)\frac{\pi}{n}.
\]

\item A particle has velocity

\[
v(t)=t^2+1
\]

meters per second for \(0\le t\le3\).  Write a right endpoint sum for its displacement.

The displacement is

\[
\int_0^3 (t^2+1)\,dt.
\]

With

\[
\Delta t=\frac3n,
\qquad
t_j=\frac{3j}{n},
\]

the sum is

\[
\sum_{j=1}^{n}
\left[
\left(\frac{3j}{n}\right)^2+1
\right]\frac3n.
\]
\end{enumerate}

\paragraph{Common traps.}
A Riemann sum fails for predictable reasons.

\[
\begin{array}{c|c}
\textbf{trap} & \textbf{fix} \\ \hline
\text{forgetting the factor }\Delta x
&
\text{every rectangle needs width}
\\[1mm]
\text{using }j\text{ instead of }x_j^*
&
\text{the function is evaluated at an }x\text{-value}
\\[1mm]
\text{using the wrong interval length}
&
\Delta x=(b-a)/n
\\[1mm]
\text{mixing left and right endpoints unintentionally}
&
\text{write }x_{j-1},x_j,\text{ or }x_j^*\text{ clearly}
\\[1mm]
\text{forgetting units}
&
\text{height units}\cdot\text{width units}=\text{accumulated units}
\end{array}
\]

\subsection{Project: estimate energy use from power data}
\label{subsec:project-estimate-energy-power-data}

Power is the rate at which energy is used.  If power is measured in kilowatts and time
is measured in hours, then

\[
\text{kilowatts}\cdot\text{hours}=\text{kilowatt-hours}.
\]

So the energy used from time \(a\) to time \(b\) is

\[
E=\int_a^b P(t)\,dt,
\]

where \(P(t)\) is power.

Suppose a small building has the following power readings during an \(8\)-hour period.

\[
\begin{array}{c|ccccccccc}
t\text{ (hours)} & 0 & 1 & 2 & 3 & 4 & 5 & 6 & 7 & 8 \\ \hline
P(t)\text{ (kW)} & 1.1 & 1.3 & 2.4 & 3.1 & 2.7 & 2.2 & 3.6 & 4.0 & 2.5
\end{array}
\]

\begin{figure}[h]
\centering
\begin{tikzpicture}
\begin{axis}[
    width=0.82\textwidth,
    height=0.45\textwidth,
    xmin=0, xmax=8.3,
    ymin=0, ymax=4.5,
    axis lines=left,
    xlabel={\(t\) (hours)},
    ylabel={\(P(t)\) (kW)},
    xtick={0,1,2,3,4,5,6,7,8},
    ytick={0,1,2,3,4},
    grid=both,
    grid style={line width=.1pt, draw=gray!25},
]
\addplot[only marks, mark=*] coordinates {
    (0,1.1) (1,1.3) (2,2.4) (3,3.1) (4,2.7)
    (5,2.2) (6,3.6) (7,4.0) (8,2.5)
};
\addplot[thick] coordinates {
    (0,1.1) (1,1.3) (2,2.4) (3,3.1) (4,2.7)
    (5,2.2) (6,3.6) (7,4.0) (8,2.5)
};
\node[anchor=west] at (axis cs:1.0,4.15) {sampled power data};
\end{axis}
\end{tikzpicture}
\caption{Energy use is accumulated power.  The area under the power graph has units of kilowatt-hours.}
\label{fig:power-data-energy-project}
\end{figure}

\paragraph{Project tasks.}
Use the data to estimate the total energy used.

\begin{enumerate}
\item Compute the left endpoint estimate

\[
L_8=
1\bigl(P(0)+P(1)+\cdots+P(7)\bigr).
\]

What are its units?

\item Compute the right endpoint estimate

\[
R_8=
1\bigl(P(1)+P(2)+\cdots+P(8)\bigr).
\]

\item Compute the trapezoid estimate

\[
T_8=\frac{L_8+R_8}{2}.
\]

Equivalently,

\[
T_8=
\frac12P(0)+P(1)+P(2)+\cdots+P(7)+\frac12P(8),
\]

because the time width is \(1\) hour.

\item Which estimate would you trust most from this data: \(L_8\), \(R_8\), or \(T_8\)?
Explain.

\item The power rises and falls during the interval.  Why is it not correct to say that
the left estimate is automatically an underestimate or automatically an overestimate?

\item Suppose the utility company charges \(\$0.18\) per kilowatt-hour.  Use your
energy estimate to estimate the cost of this \(8\)-hour period.

\item What additional data would improve the estimate?
\end{enumerate}

\begin{quote}
\textbf{What the project tests.}

An integral is not only an area under a drawn curve.  It is accumulated rate.
Power accumulates into energy, just as velocity accumulates into displacement.
\end{quote}

\subsection{Computing: compare left, right, midpoint, and trapezoid estimates}
\label{subsec:computing-compare-riemann-estimates}

Different rectangle choices can converge to the same integral at different speeds.

For an equal partition of \([a,b]\), let

\[
\Delta x=\frac{b-a}{n}.
\]

The standard estimates are:

\[
L_n=\sum_{j=0}^{n-1} f(a+j\Delta x)\Delta x,
\]

\[
R_n=\sum_{j=1}^{n} f(a+j\Delta x)\Delta x,
\]

\[
M_n=\sum_{j=1}^{n}
f\left(a+\left(j-\frac12\right)\Delta x\right)\Delta x,
\]

and

\[
T_n=
\frac{\Delta x}{2}
\left[
f(a)+2\sum_{j=1}^{n-1}f(a+j\Delta x)+f(b)
\right].
\]

The estimate \(T_n\) is the trapezoid estimate.  It replaces rectangles by trapezoids.
For equal partitions,

\[
T_n=\frac{L_n+R_n}{2}.
\]

\paragraph{Experiment.}
Use

\[
f(x)=1+x^2
\]

on \([0,2]\).  From the worked example,

\[
\int_0^2(1+x^2)\,dx=\frac{14}{3}\approx4.666667.
\]

The following table compares the estimates.

\[
\begin{array}{c|cccc}
n & L_n & R_n & M_n & T_n \\ \hline
4  & 3.750000 & 5.750000 & 4.625000 & 4.750000 \\
8  & 4.187500 & 5.187500 & 4.656250 & 4.687500 \\
16 & 4.421875 & 4.921875 & 4.664063 & 4.671875 \\
32 & 4.542969 & 4.792969 & 4.666016 & 4.667969
\end{array}
\]

The left and right estimates move toward the integral.  The midpoint and trapezoid
estimates move toward it faster in this example.

\paragraph{Computing tasks.}
Use a calculator, spreadsheet, or short program.

\begin{enumerate}
\item Extend the table to \(n=64\) and \(n=128\).

\item Compute the absolute errors

\[
|L_n-\tfrac{14}{3}|,
\qquad
|R_n-\tfrac{14}{3}|,
\qquad
|M_n-\tfrac{14}{3}|,
\qquad
|T_n-\tfrac{14}{3}|.
\]

\item What happens to the errors when \(n\) doubles?

\item Repeat the experiment for

\[
f(x)=\sqrt{1+x^2}
\]

on \([0,2]\).  Since we do not yet have an exact value, use a very fine estimate,
such as \(M_{10000}\), as a reference value.

\item Try a function that changes direction, such as

\[
f(x)=2+\sin x
\]

on \([0,2\pi]\).  Does the left estimate always stay below the right estimate?
Explain.
\end{enumerate}

This experiment prepares a later topic: numerical integration.  The integral is exact,
but in many real problems we can only approximate it.  Then the real question is not
only

\[
\text{What estimate did we compute?}
\]

but also

\[
\text{How trustworthy is the estimate?}
\]

\paragraph{Concept checks.}
Answer each question in words before using formulas.

\begin{enumerate}
\item What does \(\Delta x\) represent in a Riemann sum?

\item Why does

\[
\int_a^b f(x)\,dx
\]

measure signed area rather than ordinary area?

\item What is the difference between

\[
\int_a^b v(t)\,dt
\]

and

\[
\int_a^b |v(t)|\,dt?
\]

\item Why does changing a function at one point not change its Riemann integral?

\item Give an example of an integrable function with a jump discontinuity.

\item Give an example of a bounded function that is not Riemann integrable.

\item If \(f(x)\le g(x)\) on \([a,b]\), with \(a<b\), why must

\[
\int_a^b f(x)\,dx\le \int_a^b g(x)\,dx?
\]

\item What does the average value

\[
\frac1{b-a}\int_a^b f(x)\,dx
\]

mean geometrically?
\end{enumerate}

\subsection{Hook: an accumulation function has a slope}
\label{subsec:hook-accumulation-function-slope}

The integral adds up a function over an interval.  But what happens if one endpoint
moves?

Let

\[
A(x)=\int_0^x f(t)\,dt.
\]

Now \(A\) is a function.  Its input is the endpoint \(x\).  Its output is accumulated
signed area from \(0\) to \(x\).

\begin{figure}[h]
\centering
\begin{tikzpicture}
\begin{axis}[
    width=0.82\textwidth,
    height=0.45\textwidth,
    xmin=0, xmax=3.6,
    ymin=0, ymax=3.4,
    axis lines=left,
    xlabel={\(t\)},
    ylabel={\(f(t)\)},
    xtick={0,2,2.45,3},
    xticklabels={\(0\),\(x\),\(x+h\),\(3\)},
    ytick={0,1,2,3},
    grid=both,
    grid style={line width=.1pt, draw=gray!25},
]
\addplot[domain=0:2, samples=100, fill=gray!15, draw=none]
    {1+0.25*x+0.15*x^2} \closedcycle;
\addplot[domain=2:2.45, samples=100, fill=gray!35, draw=none]
    {1+0.25*x+0.15*x^2} \closedcycle;
\addplot[very thick, domain=0:3.2, samples=150] {1+0.25*x+0.15*x^2};
\draw[dashed] (axis cs:2,0) -- (axis cs:2,2.1);
\draw[dashed] (axis cs:2.45,0) -- (axis cs:2.45,2.51);
\node[anchor=west] at (axis cs:0.55,0.55) {\(A(x)=\int_0^x f(t)\,dt\)};
\node[rotate=90, anchor=center] at (axis cs:2.225,1.25) {new thin strip};
\node[anchor=west] at (axis cs:2.05,2.8) {\(y=f(t)\)};
\end{axis}
\end{tikzpicture}
\caption{When the endpoint moves from \(x\) to \(x+h\), the accumulation changes by the area of a thin strip.}
\label{fig:accumulation-moving-endpoint-hook}
\end{figure}

If the endpoint moves from \(x\) to \(x+h\), the change in accumulated area is

\[
A(x+h)-A(x)
=
\int_x^{x+h} f(t)\,dt.
\]

For small \(h\), the graph has little time to change.  The thin strip is approximately a
rectangle with width \(h\) and height \(f(x)\):

\[
A(x+h)-A(x)\approx f(x)h.
\]

Divide by \(h\):

\[
\frac{A(x+h)-A(x)}{h}\approx f(x).
\]

Now let \(h\to0\).  The approximation should become exact:

\[
A'(x)=f(x).
\]

This is the next great surprise.

The derivative of accumulated area is the original height.

The next chapter proves this carefully.  It is the first half of the Fundamental Theorem
of Calculus, and it completes the loop that began with the dashboard:

\[
\text{rate accumulates into amount},
\qquad
\text{and accumulated amount differentiates back to rate.}
\]